% ============================================================
% 30_flesh_twist.tex —— 骨血层 · 主链 S2（扭化）
% 底本：Kapustin–Witten, hep-th/0604151
%        §3（印刷页 16–37，source\text\p017–p038）
%        §5.1–5.2（印刷页 58–70，source\text\p059–p071）
% provenance：A = 原文明确写出；B = 原文省略或一句带过、我们重建；
%             C = 解经者判断（写作「我们判断／我们推测」）。
% 锚点格式：\anchor{src p.N} / \anchor{eq x.y} / \anchor{§x.y}，页码一律用印刷页。
% ============================================================

\section{扭化：从四维超对称杨--米尔斯到一族拓扑场论}

\noindent
\textbf{结论先行。} 本节是主链 S2 的全部内容：把四维 $\mathcal N=4$ 超杨--米尔斯（下称 $\mathcal N=4$ SYM）变成四维拓扑场论（TFT）。答案不是「找到一个拓扑超荷」，而是三层结构。第一层是表示论：用 $\mathrm{Spin}(4)_{\rm space}\times\mathrm{Spin}(6)_R$ 中一个与 $\mathrm{Spin}(4)$ 同构的\textbf{对角}子群 $\mathrm{Spin}'(4)$ 替换时空旋转群，十六个超荷里恰好有两个在该群下不变，因而变成标量。第二层是参数：这两个不变超荷的任意非零线性组合都给出一个幂零超荷，整体归一化无关，于是拓扑对称由 $\mathbb{CP}^1$ 参数化；每一个点给出一个不同的 TFT，所以得到的是\textbf{一族}而不是一个。第三层是等价性：整族的闭流形配分函数相同（vanishing theorems 保证解空间不随参数变），但算符与线算子随参数跳变；真正内禀的参数是 canonical parameter $\Psi$，$\tau$ 与 $t$ 只是它的两个坐标。二维化之后，同一套结构变成二维 sigma 模型的扭化 TFT 族，其统一语言是 Hitchin 的广义复几何\pA\anchor{src p.16, p.18, p.35, p.65}。

本节在全书中的位置：上游是骨架层的 S1（电磁对偶群同构，解释为什么 S-对偶能作用在 $\tau$ 上）；下游是 S3（§4 把四维 TFT 紧化到二维、得到 Hitchin 系统）与 S5（§6 把 Wilson--'t Hooft 线算子翻译成膜算子）。如果本节没有把「$\mathbb{CP}^1$ 族 + $\Psi$」这件事讲透，S4 与 S5 里「S-对偶在二维变成镜像对称」就会变成一句无来由的断言。全文使用节点编号 N3（= 原文 §3）、N5（= 原文 §5）、里程碑 M1（扭化后整个 $\mathbb{CP}^1$ 参数族保持拓扑性、$t=\pm\ii$ 退化）与 M2（四维 S-对偶在二维变成镜像对称）\pC\anchor{§3, §5}。

本节要回答的五个问题：
\begin{enumerate}[label=(\roman*)]
\item 为什么原始的 $\mathcal N=4$ 不能直接拿来做拓扑场论？障碍究竟在哪里？
\item 「扭化」在超荷的表示上到底做了什么？为什么恰好剩下两个标量超荷？
\item 为什么结果是一族 TFT？参数 $t$ 的地位是什么，为什么不能把它消掉？
\item 拓扑性（与度规、与 $t$ 无关）是被什么定理保证的？拓扑作用量长什么样？
\item 为什么 canonical parameter $\Psi$ 比 $\tau$ 更本质？它在二维化之后变成什么？
\end{enumerate}

\paragraph{阅读约定.}
全文在 Euclidean signature 下工作（原文指出这是扭化最自然的号差，Lorentz 号差下这套手续很别扭；当 $M=\mathbb R\times W$ 或 $\mathbb R^2\times C$ 时可以把结果继续回 Lorentz 号差）\pA\anchor{src p.16}。$M$ 默认是紧致、可定向、无边的四维流形，$G$ 是紧致单连通单李群。行文中 \pA\ 表示原文明确写出、紧随的 \anchor{} 给出印刷页或公式号；\pB\ 表示原文只说了一句或直接跳过、由我们补全，并注明重建依据；\pC\ 表示我们的判断，原文没有这样讲。本文件不重复 §1–§2 的记号。

\subsection{符号约定}

本文件新引入的符号如下（其余沿用 §1、§2；页码为印刷页，公式号为原文编号）。

{\small
\begin{longtable}{@{}p{2.45cm}p{9.55cm}p{3.35cm}@{}}
\toprule
\tabhead{符号} & \tabhead{含义} & \tabhead{出处}\\
\midrule
\endhead
$S^{\pm}$, $\hat\Gamma_E$ & Euclidean 四维 Weyl 旋量；$\hat\Gamma_E$ 在 $S^{\pm}$ 上本征值 $\mp\ii$。超对称参数取在 $S^{+}$ 上时 $\hat\Gamma_E\varepsilon=-\ii\varepsilon$ & p.17, (3.1)\\
$(2,1)$, $(1,2)$ & $\mathrm{Spin}(4)=\mathrm{SU}(2)_\ell\times\mathrm{SU}(2)_r$ 的两个双旋量表示（即 $S^{+}$、$S^{-}$） & p.17\\
$\mathrm{Spin}'(4)$ & $(1\times\kappa)\bigl(\mathrm{Spin}(4)\bigr)\subset \mathrm{Spin}(4)_{\rm space}\times\mathrm{Spin}(6)_R$，扭化所用的对角子群 & p.18, (3.2)\\
$K$, $\mathcal K$ & $\mathrm{SU}(4)_R$ 中与嵌入对易的额外 $\mathrm U(1)$ 的生成元矩阵；$\mathcal K=-\ii K$ 是其荷，等价地 $\mathcal K=-\ii\Gamma_{89}$ & p.19, (3.3), (3.11)\\
$\varepsilon_\ell$, $\varepsilon_r$ & 两个 $\mathrm{Spin}'(4)$-不变超荷参数，$\hat\Gamma_E\varepsilon_\ell=-\varepsilon_\ell$，$\hat\Gamma_E\varepsilon_r=\varepsilon_r$ & p.20, (3.7)\\
$N$ & $N=\tfrac14\sum_{\mu=0}^{3}\Gamma_{\mu+4}\Gamma_\mu$，把 $\varepsilon_\ell$ 与 $\varepsilon_r$ 配对，$N^2=-1$ & p.20, (3.8)\\
$u,v$; $t$ & $\varepsilon=u\varepsilon_\ell+v\varepsilon_r$ 的齐次坐标；仿射参数 $t=v/u$ & p.21, (3.12)\\
$Q$, $Q_\ell$, $Q_r$ & 超荷；$\delta_T$ 由 $\varepsilon$ 生成，$\delta_\ell,\delta_r$ 分别由 $\varepsilon_\ell,\varepsilon_r$ 生成 & p.22, (3.15)–(3.17)\\
$\delta_t$ & $\delta_t=\delta_\ell+t\,\delta_r$，固定 $t$ 后用来拼拓扑作用量的微分 & p.33, (3.41)\\
$A_\mu$, $\varphi_\mu$, $\sigma$ & 扭化后的玻色场：规范场、第二个 adjoint 值一形式、复标量 $\sigma=(\varphi_4-\ii\varphi_5)/\sqrt2$ & p.19\\
$\psi_\mu,\tilde\psi_\mu,\chi^{\pm},\eta,\tilde\eta$ & 费米场，$\mathcal K=1$ 的 $(2,2)$ 两份、$\mathcal K=-1$ 的 $(3,1)/(1,3)$ 二形式与两个零形式 & p.23, (3.18)\\
$P$, $H$ & 辅助场：李代数值零形式 $P$、二形式 $H$ & p.32, p.33\\
$\mathcal V_\pm(t)$, $\mathcal V_0$ & $\mathcal V_+(t)=(F-\varphi\wedge\varphi+tD\varphi)^+$，$\mathcal V_-(t)=(F-\varphi\wedge\varphi-t^{-1}D\varphi)^-$，$\mathcal V_0=D^*\varphi$ & p.28\\
$\tau$, $\theta$, $e^2$ & $\tau=\theta/2\pi+4\pi\ii/e^2$；$\mathrm{Re}\,\tau=\theta/2\pi$，$\mathrm{Im}\,\tau=4\pi/e^2$ & p.35\\
$\Psi$ & canonical parameter，$\Psi=\tfrac{\tau+\bar\tau}{2}+\tfrac{\tau-\bar\tau}{2}\tfrac{t-t^{-1}}{t+t^{-1}}$ & p.35, (3.50)\\
$n_{\mathfrak g}$ & $S$-对偶中 $\tau\to-1/(n_{\mathfrak g}\tau)$ 的归一化常数 & p.37\\
$J_\pm$, $I_w$ & 目标流形上的复结构对；超凯勒流形上以 $w\in\mathbb{CP}^1_h$ 参数化的复结构族 & p.59–60, (5.7)\\
$y$, $z$ & $\Sigma$ 与 $C$ 的复坐标（$y$ 在 $\Sigma$ 上，$z$ 在 $C$ 上） & p.61\\
$w_\pm$, $q$ & 二维族的两个 $\mathbb{CP}^1_h$ 坐标与不变量 $q=w_+/w_-$，$q=-t^2$ & p.62, (5.20)\\
$\mathcal I$, $\mathcal M(B_0)$ & $T\oplus T^*$ 上的广义复结构；闭二形式 $B_0$ 给出的 $B$-场变换 & p.65–66, (5.25)\\
$\mathcal L_\sigma$ & 由 $\sigma$ 生成的规范变换的一阶变（原文用 $\pounds_\sigma$） & p.31, (3.35)\\
\bottomrule
\end{longtable}
}

\subsection{第一个障碍：$\mathcal N=4$ 的超荷全是旋量}

\paragraph{动机.}
拓扑场论要的不是「又一个超对称理论」，而是一台算同调不变量的机器：一个\textbf{标量}、\textbf{幂零}、且定义不依赖度规的对称性 $Q$，使得可观测量取 $Q$-上同调类，配分函数与 $M$ 的度规无关。这个要求比「存在超对称」强得多。逐条看原始 $\mathcal N=4$ 的超荷，会看到两处硬障碍\pC。

\begin{enumerate}[label=(\arabic*)]
\item \textbf{超荷参数是旋量，不是标量。} 超对称变分写成 $\delta_\varepsilon=\varepsilon^{\alpha i}Q_{\alpha i}+\bar\varepsilon_{\alpha i}\bar Q^{\alpha i}$，参数 $\varepsilon$ 在 $\mathrm{Spin}(4)$ 下按 $(2,1)$ 或 $(1,2)$ 变换，没有单态分量。于是「$\delta_\varepsilon$ 作用为零」是一个依赖标架选择的陈述：在一般弯曲 $M$ 上没有协变常旋量，$\delta_\varepsilon$ 无法整体定义，「与度规无关」更无从谈起。
\item \textbf{单个超荷不幂零。} 即使 $M$ 平坦，单个超荷也满足 $\{Q,Q\}\propto \Gamma^\mu P_\mu+(\text{规范变换})$，右端是平移，不是零。于是「$Q$-上同调」没有同调代数意义：$Q$ 不是微分。
\end{enumerate}

原文的扭化定义只说「用 $\mathrm{Spin}(4)\times\mathrm{Spin}(6)$ 中另一个同构于 $\mathrm{Spin}(4)$ 的子群 $\mathrm{Spin}'(4)$ 替换 $\mathrm{Spin}(4)$」\pA\anchor{src p.18}，并没有解释为什么非这样不可。我们的重建（\pB）是：上面两条障碍恰好可以\textbf{翻译成一个表示论条件}——在 $\mathrm{Spin}(4)_{\rm space}\times\mathrm{Spin}(6)_R$ 的某个对角子群下，十六个超荷里要有单态。第一个障碍由此解决（$Q$ 成了标量），第二个障碍则要再看一步超代数。

\paragraph{推导.}

\textbf{第一步：超荷的表示内容。} Euclidean 四维的 $\mathcal N=4$ 超荷共十六个分量。由原文 (3.5)(3.6) 反读，它们在 $\mathrm{Spin}(4)_{\rm space}\times\mathrm{SU}(4)_R$ 下变换为\pB\anchor{src p.19–20}
\begin{equation}
\underbrace{\mathbf{(2,1)}\otimes\bar{\mathbf 4}}_{\text{左手 }8\text{ 个}}\;\oplus\;\underbrace{\mathbf{(1,2)}\otimes\mathbf 4}_{\text{右手 }8\text{ 个}},
\qquad \dim=2\cdot 4+2\cdot 4=16 .
\label{eq:sixteen}
\end{equation}
其中 $\mathrm{Spin}(4)_{\rm space}$ 只作用在第一个因子上，$\mathrm{SU}(4)_R\cong\mathrm{Spin}(6)_R$ 作用在第二个因子上。关键事实：$\mathrm{Spin}(4)$ 的这两个双旋量表示都是二维的、没有一维分量，所以\textbf{无论怎么组合，单个超荷都是旋量}；任何「标量超荷」都必须来自旋量指标与 R-对称指标的对角配对\pB。

\textbf{第二步：幂零性与「不变算符」论证。} 把超代数写成（示意形式，略去中心荷）\pB
\begin{equation}
\{Q_{\alpha i},Q_{\beta j}\}=\delta_{ij}(\Gamma^\mu)_{\alpha\beta}P_\mu+\cdots,
\label{eq:susyalg}
\end{equation}
取 $Q=\varepsilon^{\alpha i}Q_{\alpha i}$，则
\begin{equation}
Q^2=\tfrac12\,\varepsilon^{\alpha i}\varepsilon^{\beta j}\{Q_{\alpha i},Q_{\beta j}\}
=c^\mu P_\mu+(\text{规范变换}),
\qquad
c^\mu=\tfrac12\,\varepsilon^{\alpha i}\varepsilon^{\beta j}\delta_{ij}(\Gamma^\mu)_{\alpha\beta}.
\label{eq:Qsq}
\end{equation}
未被扭化时，$\varepsilon$ 在原来的 $\mathrm{Spin}(4)$ 下不是单态，所以 $c^\mu$ 完全可以非零（对一个 Weyl 旋量，$c^\mu$ 还是类光矢量）；这就是障碍 (2) 的定量形式\pB。

\textbf{第三步：扭化如何同时修好两件事。} 若 $\varepsilon$ 在某个新群 $\mathrm{Spin}'(4)$ 下是单态，那么 $Q$ 是标量（解决障碍 (1)）；同时 $Q$ 不变 $\Rightarrow$ $Q^2=[Q,Q\}$ 作为算符也不变，于是 $Q^2$ 的右端只能是 $\mathrm{Spin}'(4)$-不变算符。而 (2.15) 的右端只有两类东西：平移 $P_\mu$（在 $\mathbf 4_v$ 里）与规范变换（不变）。$\mathbf 4_v$ 中没有不变矢量，所以 $P_\mu$ 的系数 $c^\mu$ 必须为零，即
\begin{equation}
c^\mu=0\quad\Longrightarrow\quad Q^2=0\ \ (\text{模规范变换}).
\label{eq:Qnilpotent}
\end{equation}
这正是原文「(2.15) 的右端没有 $\mathrm{Spin}'(4)$-不变算符，所以 $Q^2$ 必为零」的论证\pA\anchor{src p.18}。我们补的一点（\pB）是把「没有不变算符」拆成两步：先指出 $Q$ 不变迫使 $Q^2$ 不变，再指出唯一可能破坏幂零性的算符是平移、而它在 $\mathbf 4_v$ 里。

\paragraph{结论的作用.}
这一小节把构造 TFT 的问题\textbf{改造成了纯表示论问题}：找同态 $\kappa:\mathrm{Spin}(4)\to\mathrm{Spin}(6)$，使得旋量表示里出现单态。原文 §3.1 余下的全部内容，就是这一句话的展开\pA\anchor{src p.18–22}。它还解释了后面一切结果的形状：既然 $Q$ 只是「不变子空间里的一个方向」，而这个不变子空间在 $\mathcal N=4$ 里恰好是二维复的，那么\textbf{参数族的出现就不是偶然，而是必然}——这正是下一小节的主题。同一个论证还给出一个以后要反复用的精确形式：$\delta_t^2=-\ii(1+t^2)\mathcal L_\sigma$，即 $Q^2$ 至多是一个规范变换，$t=\pm\ii$ 时才严格为零\pA\anchor{src p.33, (3.42)}。

\begin{warn}
\textbf{排除错误直觉（一）：扭化把 $\mathcal N=4$ 变简单了。} 错。在 $\mathbb R^4$ 上，扭化前后的\textbf{作用量、场内容、超代数完全相同}；变的只是「我们把哪个子代数叫作超对称、把 $\mathrm{SU}(4)_R$ 的哪一部分与时空旋转等同起来」。所谓扭化理论，是同一个理论在另一组变量下的\textbf{上同调截断}：只保留 $Q$-闭算符、把 $Q$-恰当算符视为零。真正被改变的是\textbf{我们提什么问题}，而不是理论本身。反过来，只有扭化之后理论才能定义在一般弯曲 $M$ 上——未扭化的 $\mathcal N=4$ SYM 并不存在于一般四维流形\pC。
\end{warn}

\begin{stuck}
为什么必须在 Euclidean signature 下做？原文只说「Euclidean 最自然、Lorentz 下很别扭」\pA\anchor{src p.16}。我们的理解（\pB）：Lorentz 群 $\mathrm{Spin}(3,1)\cong\mathrm{SL}(2,\mathbb C)$ 的旋量表示是复表示，$\mathrm{SU}(4)_R=\mathrm{Spin}(6)$ 的 $\mathbf 4,\bar{\mathbf 4}$ 也是复表示，两边的「同型配对」要带复共轭，扭化所依赖的对角嵌入会牵扯到实结构；而 Euclidean 下 $S^{\pm}$ 是 pseudoreal 的，$(2,1)$ 与 $\mathrm{SU}(2)$ 的表示才有干净的对应\pC。这就是全节都在 Euclidean 号差下工作的原因。
\end{stuck}

\subsection{$\mathrm{SU}(4)_R\cong\mathrm{Spin}(6)$ 与超荷的重排}

\paragraph{动机.}
现在要具体找出 $\kappa$。要判断「哪些超荷不变」，必须把十六个超荷在扭化前后的量子数都写清楚：扭化前按 $\mathrm{Spin}(4)_{\rm space}\times\mathrm{SU}(4)_R$ 分类，扭化后按 $\mathrm{Spin}'(4)\times\mathrm U(1)_{\mathcal K}$ 分类。这一步在原文里是两行分解式 (3.5)(3.6)，但它其实是整个 §3 的枢纽，值得把每一步的表示论依据写全。

\paragraph{推导.}

\textbf{第一步：R-对称的几何来源。} $\mathcal N=4$ 的六个标量 $\varphi_0,\dots,\varphi_5$ 在 $\mathrm{SO}(6)_R\cong\mathrm{SU}(4)_R$ 的 $\mathbf 6_v$ 里，超荷带 $\mathbf 4$ 或 $\bar{\mathbf 4}$。GL twist 选择的 $\kappa$ 把 $\mathrm{Spin}(4)=\mathrm{SU}(2)_\ell\times\mathrm{SU}(2)_r$ \textbf{分块对角}地嵌入 $\mathrm{SU}(4)_R$，使得\pA\anchor{src p.19, (3.2)}
\begin{equation}
\mathbf 4=(\mathbf 2,\mathbf 1)_{1}\oplus(\mathbf 1,\mathbf 2)_{-1},
\qquad
\bar{\mathbf 4}=(\mathbf 2,\mathbf 1)_{-1}\oplus(\mathbf 1,\mathbf 2)_{1},
\label{eq:fourbar}
\end{equation}
其中下标是 $\mathcal K$ 荷（定义见下）。这个嵌入还与一个额外的 $\mathrm U(1)$ 对易，其生成元矩阵是 $2\times2$ 分块的 $K=\ii\,\mathrm{diag}(1,-1)$\pA\anchor{src p.19, (3.3)}。严格地说，整体群是 $(\mathrm{SU}(2)_\ell\times\mathrm{SU}(2)_r\times\mathrm U(1))/\mathbb Z_2$；这个 $\mathrm U(1)$ 是 $\mathrm{SO}(4)\times\mathrm U(1)\subset \mathrm{SO}(6)$ 中 $\mathrm U(1)$ 的双重覆盖\pA\anchor{src p.19, 脚注 7–8}。荷的约定是：$\mathcal K=-\ii K$，所以 (3.5)(3.6) 里的指数就是 $\mathcal K$ 本征值。

\textbf{第二步：SO(6) 语言——六个标量的重排。} 在 $\mathrm{SO}(6)$ 语言里
\begin{equation}
\mathbf 6_v=\Lambda^2\mathbf 4=(\mathbf 2,\mathbf 2)_0\oplus(\mathbf 1,\mathbf 1)_{2}\oplus(\mathbf 1,\mathbf 1)_{-2},
\label{eq:sixv}
\end{equation}
即 $\mathrm{SO}(4)\times\mathrm{SO}(2)\hookrightarrow\mathrm{SO}(6)$ 的分块嵌入\pA\anchor{src p.19, (3.4)}。于是四个标量 $\varphi_0,\dots,\varphi_3$ 组成 $\mathbf 4_v$，在新的旋转群下\textbf{升格为一个一形式} $\varphi_\mu$；剩下两个 $\varphi_4,\varphi_5$ 被 $\mathrm{SO}(2)$ 旋转，组合成
\begin{equation}
\sigma=\frac{\varphi_4-\ii\varphi_5}{\sqrt2}\ (\mathcal K=2),
\qquad
\bar\sigma=\frac{\varphi_4+\ii\varphi_5}{\sqrt2}\ (\mathcal K=-2).
\label{eq:sigma}
\end{equation}
\textbf{结论：扭化后的玻色场是「规范场 $A_\mu$ $+$ 第二个 adjoint 一形式 $\varphi_\mu$ $+$ 复标量 $\sigma,\bar\sigma$」。}这个额外的 $\varphi_\mu$ 是全部物理的来源：没有它，方程 (3.29) 就退化成纯 Yang--Mills；有了它，固定点方程才成为 Hitchin 方程那一类东西\pB\anchor{src p.19}。

\textbf{第三步：超荷的重排（本小节的核心计算）。} 把 (3.5)(3.6) 的每一步都写出来。用 $\mathrm{SU}(2)$ 的 $2\otimes2=1\oplus3$、$2\otimes1=2$，以及「$\mathcal K$ 荷相加」，得到\pA\anchor{src p.20, (3.5)–(3.6)}
\begin{align}
\mathbf{(2,1)}^{0}\otimes\Bigl[\mathbf{(2,1)}^{-1}\oplus\mathbf{(1,2)}^{+1}\Bigr]
&=\underbrace{\bigl[\mathbf{(2,1)}\otimes\mathbf{(2,1)}\bigr]^{-1}}_{\mathbf{(1,1)}^{-1}\oplus\mathbf{(3,1)}^{-1}}
\;\oplus\;
\underbrace{\bigl[\mathbf{(2,1)}\otimes\mathbf{(1,2)}\bigr]^{+1}}_{\mathbf{(2,2)}^{+1}},
\label{eq:branch5}\\[2pt]
\mathbf{(1,2)}^{0}\otimes\Bigl[\mathbf{(2,1)}^{+1}\oplus\mathbf{(1,2)}^{-1}\Bigr]
&=\underbrace{\bigl[\mathbf{(1,2)}\otimes\mathbf{(2,1)}\bigr]^{+1}}_{\mathbf{(2,2)}^{+1}}
\;\oplus\;
\underbrace{\bigl[\mathbf{(1,2)}\otimes\mathbf{(1,2)}\bigr]^{-1}}_{\mathbf{(1,1)}^{-1}\oplus\mathbf{(1,3)}^{-1}}.
\label{eq:branch6}
\end{align}
推导依据只有两条：（i）左边的整体 $\mathcal K$ 荷（旋量部分取 $0$，$\mathbf 4$ 与 $\bar{\mathbf 4}$ 的分解由 \eqref{eq:fourbar} 给出）相加；（ii）两个 $\mathrm{SU}(2)$ 因子各自做张量积，且 $\mathrm{Spin}'(4)$ 的表示由\textbf{对角}作用合成，所以 $\mathbf{(2,1)}\otimes\mathbf{(2,1)}=(\mathbf{1},\mathbf{1})\oplus(\mathbf{3},\mathbf{1})$ 这一类等式成立\pB。

\textbf{第四步：数单态。} 在 (3.5)(3.6) 的右端里，只有 $\mathbf{(1,1)}$ 在 $\mathrm{Spin}'(4)$ 下不变，而每个分解恰有一个 $\mathbf{(1,1)}$，并且两个的 $\mathcal K$ 都是 $-1$。所以\textbf{恰好两个不变超荷参数}：$\varepsilon_\ell$（在 (3.5) 里，是四维意义下的左手旋量）与 $\varepsilon_r$（在 (3.6) 里，右手）\pA\anchor{src p.20}。两者由
\begin{equation}
\varepsilon_r=N\varepsilon_\ell,\qquad
N=\tfrac14\sum_{\mu=0}^{3}\Gamma_{\mu+4}\Gamma_\mu,\qquad
N^2=-1\ \ (\text{在不变旋量上})
\label{eq:N}
\end{equation}
联系起来：$N$ 与 $\mathrm{Spin}'(4)$ 对易、与 $\hat\Gamma_E$ 反对易，所以 (3.8) 作为「自然选择」是自洽的\pA\anchor{src p.20, (3.8)–(3.9)}。等价地有 $\varepsilon_\ell=-N\varepsilon_r$ 与
\begin{equation}
\Gamma_{\mu+4}\varepsilon_\ell=-\Gamma_\mu\varepsilon_r,
\qquad
\Gamma_{\mu+4}\varepsilon_r=\Gamma_\mu\varepsilon_\ell .
\label{eq:gammarels}
\end{equation}

\begin{figure}[htbp]
\centering
\begin{tikzpicture}[font=\scriptsize,
  cell/.style={draw=kwgray!55, minimum width=0.72cm, minimum height=0.72cm, inner sep=0pt},
  blk/.style={draw=kwblue, thick, rounded corners=2pt},
  lab/.style={font=\scriptsize}]
% ---------- 左面板：(3.5) ----------
\begin{scope}
  \foreach \i in {0,1}{\foreach \j in {0,1,2,3}{
    \node[cell] at (\j*0.72, -\i*0.72) {};}}
  \draw[blk] (-0.36,0.36) rectangle (1.08,-1.08);
  \draw[blk] (1.08,0.36) rectangle (2.52,-1.08);
  \node[lab,anchor=east] at (-0.46,-0.36) {$S^{+}$};
  \node[lab,anchor=south] at (0.36,0.44) {$\mathbf{(2,1)}^{-1}$};
  \node[lab,anchor=south] at (1.80,0.44) {$\mathbf{(1,2)}^{+1}$};
  \node[lab,anchor=north,align=center] at (0.36,-1.22)
    {$\mathbf{(1,1)}^{-1}$\\[1pt]$\oplus\,\mathbf{(3,1)}^{-1}$\\[1pt]$\Rightarrow\ \varepsilon_\ell$};
  \node[lab,anchor=north,align=center] at (1.80,-1.22)
    {$\mathbf{(2,2)}^{+1}$\\[1pt]无单态};
  \node[lab,anchor=south] at (1.08,1.12) {$(3.5)$：$S^{+}\otimes\bar{\mathbf 4}$};
\end{scope}
% ---------- 右面板：(3.6) ----------
\begin{scope}[shift={(5.0,0)}]
  \foreach \i in {0,1}{\foreach \j in {0,1,2,3}{
    \node[cell] at (\j*0.72, -\i*0.72) {};}}
  \draw[blk] (-0.36,0.36) rectangle (1.08,-1.08);
  \draw[blk] (1.08,0.36) rectangle (2.52,-1.08);
  \node[lab,anchor=east] at (-0.46,-0.36) {$S^{-}$};
  \node[lab,anchor=south] at (0.36,0.44) {$\mathbf{(2,1)}^{+1}$};
  \node[lab,anchor=south] at (1.80,0.44) {$\mathbf{(1,2)}^{-1}$};
  \node[lab,anchor=north,align=center] at (0.36,-1.22)
    {$\mathbf{(2,2)}^{+1}$\\[1pt]无单态};
  \node[lab,anchor=north,align=center] at (1.80,-1.22)
    {$\mathbf{(1,1)}^{-1}$\\[1pt]$\oplus\,\mathbf{(1,3)}^{-1}$\\[1pt]$\Rightarrow\ \varepsilon_r$};
  \node[lab,anchor=south] at (1.08,1.12) {$(3.6)$：$S^{-}\otimes\mathbf 4$};
  \node[lab,anchor=west,align=left] at (2.92,-0.36)
    {每块 $2\times2$：\\[1pt]单态是\textbf{对角组合}，\\[1pt]不占单个格子};
\end{scope}
\end{tikzpicture}
\caption{十六个超荷按 $\mathrm{Spin}'(4)\times\mathrm U(1)_{\mathcal K}$ 的重排（式 (3.5)(3.6) 的格点图）。每个 $2\times2$ 方块是一个 $\mathrm{SU}(2)$ 双线性；只有含 $\mathbf{(1,1)}$ 的方块里存在单态方向，即 $\varepsilon_\ell$ 与 $\varepsilon_r$。原文未画此图，此处按原文分解绘制\pB\anchor{src p.20}。}
\label{fig:grid}
\end{figure}

\begin{table}[htbp]
\centering\small
\caption{扭化后场的量子数（$\mathrm{Spin}'(4)$ 表示与 $\mathcal K$ 荷）。数据取自原文 §3.1，最后一行与 $\eta,\tilde\eta$ 同表示，故列在一起\pA\anchor{src p.23}。}
\label{tab:fields}
\begin{tabular}{@{}llcl@{}}
\toprule
\tabhead{场} & \tabhead{$\mathrm{Spin}'(4)$ 表示} & \tabhead{$\mathcal K$} & \tabhead{说明}\\
\midrule
$A_\mu$ & $(2,2)$ & $0$ & 规范场\\
$\varphi_\mu$ & $(2,2)$ & $0$ & 由 $\mathbf 4_v$ 的四个标量升格为一形式\\
$\sigma,\ \bar\sigma$ & $(1,1)$ & $+2,\,-2$ & 由两个 $\mathrm{SO}(2)$ 标量组合\\
$\psi_\mu,\ \tilde\psi_\mu$ & $(2,2)$ & $+1$ & 两份一形式费米子\\
$\chi^{+}_{\mu\nu},\ \chi^{-}_{\mu\nu}$ & $(3,1),\ (1,3)$ & $-1$ & 二形式 $\chi$ 的自对偶／反自对偶部分\\
$\eta,\ \tilde\eta$ & $(1,1)$ & $-1$ & 两个零形式费米子\\
$\varepsilon_\ell,\ \varepsilon_r$ & $(1,1)$ & $-1$ & 超荷参数（与 $\eta$ 同表示）\\
\bottomrule
\end{tabular}
\end{table}

\paragraph{结论的作用.}
这一小节交出三件东西，全部被后文消费：（i）$\varepsilon_\ell,\varepsilon_r$ 的存在，使下一小节能构造 $\mathbb{CP}^1$ 族；（ii）$\mathrm{Spin}'(4)\times\mathrm U(1)_{\mathcal K}$ 的分次，使 $\delta_T$ 成为\textbf{提高 $\mathcal K$ 一阶}的微分，从而「规范不变算符空间 $=\mathcal K$-分次复形、微分 $Q$」\pA\anchor{src p.23}；（iii）玻色场的重排（规范场 $+$ 一形式 $+$ 复标量），正是 §3.2 的方程 (3.29) 与 S3 的 Hitchin 系统所用的场内容\pB。这张表也预告了下一小节为什么 $t$ 是真参数：$\mathcal K$ 的分配不是任意的。

\begin{warn}
\textbf{排除错误直觉（二）：R-对称纯粹是内部对称，与时空无关，所以扭化不可能动它。} 错在最后一句。扭化\textbf{没有改变} R-对称（$\mathrm{SU}(4)_R$ 仍然是理论的一个对称群），改变的是「我们把哪一个子群叫作旋转群」：新的旋转群是 $\mathrm{Spin}(4)_{\rm space}\times\mathrm{Spin}(6)_R$ 的\textbf{对角}子群，它同时作用在旋量指标与 R-对称指标上。正因为两个指标被同一个 $g$ 同时旋转，一个旋量指标才能与一个 R-对称指标「配对」成单态。若把 R-对称当成与时空无关的旁观者，就永远看不出为什么 (3.5) 里会出现 $\mathbf{(1,1)}$\pB。
\end{warn}

\caveat{(3.5)(3.6) 里指数的符号是\textbf{荷相加}，不是相乘；$\mathbf 4$ 与 $\bar{\mathbf 4}$ 的指数互为相反数（见 \eqref{eq:fourbar}）。抄写时最常见的错误是把 (3.5) 里 $\bar{\mathbf 4}$ 的指数写成与 $\mathbf 4$ 相同，这样两个单态的 $\mathcal K$ 会变成 $-1$ 与 $+1$，下一小节的整个论证就崩了——见下一小节的卡住点\pA\anchor{src p.19–20}。}

\subsection{为什么是一族：两个不变超荷与 $\mathbb{CP}^1$}

\paragraph{动机.}
两个不变超荷意味着「拓扑超荷」不是唯一的。为什么它们的\textbf{任意}线性组合仍给出合法的 TFT？为什么这个组合参数消不掉？又为什么参数空间恰是 $\mathbb{CP}^1$ 而不是别的？这三个问题的答案合起来，就是「为什么是一族」的全部内容。

\paragraph{推导.}

\textbf{第一步：$\mathbb{CP}^1$ 从哪来。} 令
\begin{equation}
\varepsilon=u\,\varepsilon_\ell+v\,\varepsilon_r,
\qquad
Q=u\,Q_\ell+v\,Q_r,
\qquad
(u,v)\in\mathbb C^2\setminus\{0\},
\label{eq:epsuv}
\end{equation}
则整体标度 $(u,v)\sim\lambda(u,v)$ 不改变任何物理（它只是把同一个超荷乘一个非零数），所以拓扑对称由 $\mathbb{CP}^1$ 参数化，仿射坐标取 $t=v/u$\pA\anchor{src p.21, (3.12)}。相应的变分是 $\delta_T=\varepsilon_\ell\delta_\ell+\varepsilon_r\delta_r$，即
\begin{equation}
\delta_T\Phi=[Q,\Phi\},\qquad
\delta_\ell\Phi=[Q_\ell,\Phi\},\qquad
\delta_r\Phi=[Q_r,\Phi\}.
\label{eq:deltaT}
\end{equation}
$Q^2=0$（模规范变换）由上一小节的表示论论证保证，与 $t$ 无关。

\textbf{第二步：这个 $\mathbb{CP}^1$ 比抽象射影空间多了结构。} 原文指出，正因为 $\varepsilon_r$ 由 $\varepsilon_\ell$ \textbf{自然}确定（通过 $N$，见 (3.8)–(3.10)），$\mathbb{CP}^1$ 上还有一个额外的复结构，$t$ 才是「自然」的仿射坐标\pA\anchor{src p.22}。我们的读法（\pB）是：$N$ 给出一个复结构，把 $\varepsilon_\ell$ 与 $\varepsilon_r$ 配对；$t=0$（$\varepsilon=\varepsilon_\ell$）与 $t=\infty$（$\varepsilon=\varepsilon_r$）是两个「纯」点。这个结构以后会以「把 $t$ 送到 $1/t$ 的反定向对称」的形式再次出现（§3.5，式 (3.56)）。

\textbf{第三步：为什么 $\mathrm U(1)_{\mathcal K}$ 消不掉 $t$——全节最脆的一步。} 三个不变超荷的荷全是 $\mathcal K=-1$\pA\anchor{src p.22}，因此任一 $Q=uQ_\ell+vQ_r$ 的 $\mathcal K$ 都是 $+1$，$\delta_T$ 把 $\mathcal K$ 提高一阶。假如 $\varepsilon_\ell$ 与 $\varepsilon_r$ 的 $\mathcal K$ 不同，那么 $\mathrm U(1)_{\mathcal K}$ 会把 $(u,v)$ 相对地乘上 $e^{\ii\alpha}$ 与 $e^{-\ii\alpha}$，从而把 $t$ 变成 $e^{2\ii\alpha}t$：整个 $\mathbb{CP}^1$ 上的理论彼此等价，$t$ 就没有物理意义。正因两者 $\mathcal K$ 相同，$\mathrm U(1)_{\mathcal K}$ 只把 $(u,v)$ 整体乘一个相，在射影空间上无作用，$t$ 才是真参数\pA\anchor{src p.22}。原文对此有一句非常重的评论：正因几乎没有平凡等价，$S$-对偶带来的非平凡等价才有可能产生几何 Langlands 这样的结构\pA\anchor{src p.22}。

\textbf{第四步：哪些 $t$ 的变换是平凡的。} $\mathrm{SU}(4)_R$ 的中心 $I$ 作用为 $\varepsilon_\ell\to\ii\varepsilon_\ell$、$\varepsilon_r\to\ii^{-1}\varepsilon_r$，于是在 $\mathbb{CP}^1$ 上
\begin{equation}
I:\ t\longmapsto -t .
\label{eq:tminus}
\end{equation}
所以 $t$ 与 $-t$ 给出同一个 TFT。原文强调这是\textbf{经典作用量的对称性}（而非 $S$-对偶），并且是\textbf{唯一}的平凡等价\pA\anchor{src p.22, (3.14)}。下一步要用的 $S$-对偶作用则不是平凡的：由 R-对称对超荷参数的变换 (2.25)，$\varepsilon_\ell\to e^{\ii\varphi}\varepsilon_\ell$、$\varepsilon_r\to e^{-\ii\varphi}\varepsilon_r$，故
\begin{equation}
t\ \longmapsto\ e^{2\ii\varphi(\gamma)}\,t=\frac{c\tau+d}{|c\tau+d|}\,t,
\qquad
\gamma=\begin{pmatrix}a&b\\c&d\end{pmatrix}\in\mathrm{SL}(2,\mathbb Z),
\label{eq:taction}
\end{equation}
其中 $e^{\ii\varphi(\gamma)}=\bigl((c\tau+d)/|c\tau+d|\bigr)^{1/2}$\pA\anchor{src p.22, (3.13)}。注意 \eqref{eq:taction} 不是 $\tau$ 的全纯函数（分母带 $|\cdot|$）——这正是为什么 §3.5 需要重新定义一个参数，使 $S$-对偶变成全纯的分式线性变换。

\paragraph{结论的作用.}
一族 TFT 的骨架就此搭好：每个 $t\in\mathbb{CP}^1$ 给出一个幂零超荷 $Q_t$ 与一个（模规范变换的）上同调问题，规范不变算符空间按 $\mathcal K$ 分次，$\delta_T$ 是提高一阶的微分\pA\anchor{src p.23}。这个族的结构由两件事决定：$\mathbb{CP}^1$ 上的自然坐标 $t$（式 (3.12)）与唯一平凡等价 $t\sim-t$（式 (3.14)）。下一小节把 $t$ 的取值与方程的性质对应起来，§3.3–§3.4 再证明整族在物理上「算同一个不变量」。

\begin{warn}
\textbf{排除错误直觉（三）：$t$ 是耦合常数。} 错。$t$ 不出现在作用量里作为任何项的系数；它参数化的是「用哪个超荷」，即拓扑对称的选择\pA\anchor{src p.21}。真正的耦合是 $\tau$（$\mathrm{Im}\,\tau=4\pi/e^2$、$\mathrm{Re}\,\tau=\theta/2\pi$）。把 $t$ 当耦合会立刻在 $S$-对偶处出错：$S$-对偶\textbf{同时}作用在 $\tau$ 与 $t$ 上（\eqref{eq:taction}），并不保持 $t$ 不变。正确的图像是：$(\tau,t)$ 不是两个独立耦合，物理上只有组合 $\Psi$（§3.5）——但这恰恰是后面才证明的事，不能提前当地理常识用。
\end{warn}

\begin{warn}
\textbf{排除错误直觉（四）：$\mathbb{CP}^1$ 上的理论通过连续变形互相连通，所以都等价。} 错。$Q_t$ 连续依赖 $t$，但\textbf{上同调的维数会跳变}：$Q_t$-闭算符的集合随 $t$ 变，$t=\pm\ii$ 处结构退化（§3.2 的方程从椭圆族变成复平坦联络）。被证明与 $t$ 无关的只是\textbf{无插入的闭流形配分函数}（§3.3）；算符、边界条件与线算符都随 $t$ 变，这正是全文后面「在特定 $t$ 才存在的算子」的来源\pA\anchor{src p.27–28}。
\end{warn}

\begin{stuck}
请在这里停下来想清楚一件事：$\mathrm U(1)_{\mathcal K}$ 对 $t$ 无作用，是因为两个不变超荷的 $\mathcal K$ \textbf{相同}。这一点不是自动的，而是 $\mathcal N=4$ 的表示论恰好如此（$\varepsilon_\ell,\varepsilon_r$ 都是 (3.5)(3.6) 里唯一的 $\mathbf{(1,1)}^{-1}$）。如果读者跳过了这里的「脆性」，就会误以为参数族是「随便一个对称群破缺」都会有的现象，从而看不懂 §5 里为什么 $S$-对偶在二维能变成镜像对称\pC\anchor{src p.22}。
\end{stuck}

\subsection{原文 §3.2：固定点方程与「一族 A-模型」}

\paragraph{动机.}
有了 $Q_t$，下一步是把「拓扑」落到可以计算的方程上。路径积分在 $Q$-不变场上有分布：若 $Q$-不变玻色场空间 $Y$ 光滑、规范群自由作用，则可以用高斯近似（单圈）计算；一般情形下配分函数只由 $Y$ 的有限阶无穷小邻域决定\pA\anchor{src p.25}。所以\textbf{只需要找 $Q$-不变玻色场}，即对所有费米场解 $\delta_T\Psi=0$。这一步把「一族 TFT」翻译成「一族方程」，是本节与 S3 的接口。

\paragraph{推导.}

\textbf{第一步：写出变换律。} 由 (3.23)–(3.28)，在扭化后的场变量下\pA\anchor{src p.24, (3.23)–(3.28)}
\begin{equation}
\begin{aligned}
\delta_T A_\mu&=\ii u\,\psi_\mu+\ii v\,\tilde\psi_\mu, &
\delta_T\varphi_\mu&=\ii v\,\psi_\mu-\ii u\,\tilde\psi_\mu,\\
\delta_T\sigma&=0, &
\delta_T\bar\sigma&=\ii u\,\eta+\ii v\,\tilde\eta,\\
\delta_T\chi^{+}&=u\,(F-\varphi\wedge\varphi)^{+}+v\,(D\varphi)^{+}, &
\delta_T\chi^{-}&=v\,(F-\varphi\wedge\varphi)^{-}-u\,(D\varphi)^{-},\\
\delta_T\eta&=v\,D^*\varphi+u\,[\sigma,\bar\sigma], &
\delta_T\tilde\eta&=-u\,D^*\varphi+v\,[\sigma,\bar\sigma],\\
\delta_T\psi&=u\,D\sigma+v\,[\varphi,\sigma], &
\delta_T\tilde\psi&=v\,D\sigma-u\,[\varphi,\sigma].
\end{aligned}
\label{eq:transforms}
\end{equation}
其中 $D=\dd+[A,\cdot]$，$D^*\varphi=\star D\star\varphi=D_\mu\varphi_\mu$。$\delta_T\sigma=0$ 不是假设：$\sigma$ 的 $\mathcal K=2$，而 $\delta_T$ 把 $\mathcal K$ 提高一阶，$\mathcal K=3$ 处没有场\pA\anchor{src p.24}。

\textbf{第二步：$\chi^{\pm}$ 的固定点条件。} $\delta_T\chi^{\pm}=0$ 给出
\begin{equation}
(F-\varphi\wedge\varphi+tD\varphi)^{+}=0,
\qquad
(F-\varphi\wedge\varphi-t^{-1}D\varphi)^{-}=0,
\qquad
D^*\varphi=0,
\label{eq:fixed}
\end{equation}
前两式即 $\mathcal V_\pm(t)=0$。第三个条件来自 $\delta_T\eta=\delta_T\tilde\eta=0$ 的一个 $t$-相关线性组合\pA\anchor{src p.25, (3.29)}。把这个组合写清楚是原文省略的一步（\pB）：方程组
\begin{equation}
\begin{pmatrix}v&u\\-u&v\end{pmatrix}
\begin{pmatrix}D^*\varphi\\ [\sigma,\bar\sigma]\end{pmatrix}=0
\qquad\Longrightarrow\qquad
\det=u^2+v^2=u^2(1+t^2),
\label{eq:etaMatrix}
\end{equation}
所以在 $t\ne\pm\ii$ 时它强制 $D^*\varphi=0$ 且 $[\sigma,\bar\sigma]=0$；而在 $t=\pm\ii$ 时矩阵退化（秩 $1$），只留下一个线性组合 $D^*\varphi\pm\ii[\sigma,\bar\sigma]=0$。这正是原文所说「$t=\pm\ii$ 时细节略有不同」的精确内容\pA\anchor{src p.25–27}。

\textbf{第三步：$t\ne\pm\ii$ 时方程不依赖 $\sigma$。} 未破缺超对称还要求 (3.30)：$D\sigma+t[\varphi,\sigma]=0$、$D\bar\sigma-t^{-1}[\varphi,\bar\sigma]=0$。对 $t\ne\pm\ii$，这两式与规范固定条件一起推出 $D\sigma=[\varphi,\sigma]=0$ 与 $[\sigma,\bar\sigma]=0$，于是 $\sigma$ 生成的规范变换保持解；若超对称场是不可约的，则 $\sigma=0$\pA\anchor{src p.26, (3.30)–(3.31)}。本文只考虑不可约情形（可约情形原文指向 [26]）。所以对 $t\ne\pm\ii$，固定点方程就是关于 $(A,\varphi)$ 的
\begin{equation}
\mathcal V_+(t)=\mathcal V_-(t)=\mathcal V_0=0,
\label{eq:fixedBPS}
\end{equation}
而 $t$ 只以两种方式出现——这正是「一族方程」的字面含义。

\textbf{第四步：两个极端 $t=\pm\ii$。} 此时 $-t^{-1}=t$，令 $\mathcal A=A+\ii\varphi$（$G_{\mathbb C}$ 值联络，$\mathcal F=\dd\mathcal A+\mathcal A\wedge\mathcal A$）。直接计算给出
\begin{equation}
\mathcal A=A+\ii\varphi
\quad\Longrightarrow\quad
\mathcal F=(F-\varphi\wedge\varphi)+\ii D\varphi,
\label{eq:FcalA}
\end{equation}
所以 $\mathcal F=0$ 等价于 $(F-\varphi\wedge\varphi)^{+}=-\ii(D\varphi)^{+}$ 与 $(F-\varphi\wedge\varphi)^{-}=-\ii(D\varphi)^{-}$；在 $t=\ii$ 时 (3.29) 恰好给出这两式。$t=-\ii$ 时同样成立，但要用共轭联络 $\bar{\mathcal A}=A-\ii\varphi$\pB\anchor{src p.25}。几何上，$\mathcal F=0$ 的解就是平坦 $G_{\mathbb C}$ 联络，即同态 $\pi_1(M)\to G_{\mathbb C}$（模共轭）；再加上 $D^*\varphi=0$ 并除以 $G$-规范变换，由 Corlette 定理（二维情形是 Donaldson 的结果）等价于「去掉最后一个方程、令 $\mathcal F=0$、加稳定性条件、除以 $G_{\mathbb C}$-规范变换」，即同态模空间 $Y_M(G)$\pA\anchor{src p.25}。

\textbf{第五步：实 $t$ 与非实 $t$。} 对非实 $t$，(3.29) 对取实李代数值的 $(A,\varphi)$ 是超定的、非椭圆的；分离 Hermitian 与反 Hermitian 部分可推出 $F=0$。真正有几何意义的是实 $t$：此时方程取实值，得到一个实椭圆方程族，$t\ne0,\infty$ 处正则；$t=0$ 时第二式乘以 $t$ 给出 $(D\varphi)^{-}=0$，$t=\infty$ 时第一式乘以 $t^{-1}$ 给出 $(D\varphi)^{+}=0$，所以整个实椭圆族由 $\mathbb{RP}^1$ 参数化\pA\anchor{src p.26}。原文进一步指出：实 $t$ 的理论属于 cohomological field theory 的框架（Donaldson 理论、A-模型都是此类），玻色场 $(A,\varphi)$ 的 $\mathcal K=0$，方程就是 (3.29)，对称性只剩规范对称；无插入配分函数由（带符号地）数孤立不可约解的个数给出，符号来自单圈费米子行列式\pA\anchor{src p.27}。

\paragraph{结论的作用.}
一族方程就此得到，它同时向两个方向供料：（i）向 §3.3 提供「解空间」，使「族是否算同一个不变量」变成「解空间是否随 $t$ 变」；（ii）向 S3 提供几何对象——$t=\pm\ii$ 时的平坦 $G_{\mathbb C}$ 联络模空间，在 $M=\Sigma\times C$ 上会变成 $\MH$ 与 Hitchin 纤维化\pA\anchor{src p.25, p.30}。原文还强调，实 $t$ 与 $t=\pm\ii$ 分别「类似二维 A-模型与 B-模型」：同调类模型里超对称场是常映射，而常映射与平坦 $G_{\mathbb C}$ 联络都可以用超定方程（$\dd\Phi=0$ 与 $\mathcal F=0$）描述\pA\anchor{src p.28}。

\begin{warn}
\textbf{排除错误直觉（五）：原文说「一族 A-模型」，所以它就是二维 A-模型的参数变形。} 错。原文用的词是「类似」（similar）\pA\anchor{src p.28}：四维的 (3.29) 与二维 A-模型只是\textbf{结构上}相似——都是超定方程、都带一个复结构式的数据。真正的联系要等 S3：把 $M$ 取成 $\Sigma\times C$ 并紧化之后，二维的 A-模型／B-模型才会作为极限出现（§5.1）。把四维方程直接等同于 A-模型，会漏掉 $\varphi_\mu$ 这个一形式，也会误解 $t$ 的作用。
\end{warn}

\caveat{(3.29) 的第二式是「$-\ t^{-1}$」而不是「$+\ t^{-1}$」：$\mathcal V_-(t)=(F-\varphi\wedge\varphi-t^{-1}D\varphi)^{-}$。这个符号决定了 $t\leftrightarrow-t$ 的等价（(3.14)）与 $t\leftrightarrow t^{-1}$ 的反定向对称（(3.56)），抄错一个符号，全节的对称性都对不上。}

\subsection{原文 §3.3：vanishing theorems——族为什么是同一台机器}

\paragraph{动机.}
现在有两个不同的问题要分开：（i）配分函数是否与 $t$ 无关？（ii）是否与 $M$ 的度规无关？本节只解决 (i)，而且用的办法比「逐点比较」聪明得多：证明\textbf{在某个 $t$ 的解自动是对所有 $t$ 的解}（在 Pontryagin 类非零的情形），或者至少证明解空间在 $t$ 变动下不变。这就是 vanishing theorems。它们之所以叫这个名字，是因为最重要的推论是：Pontryagin 类非零的丛上根本没有解。

\paragraph{推导.}

\textbf{定理 1（VT1）与恒等式 (3.32)。} 设 $M$ 紧致无边、$G$-丛 $E$ 的 Pontryagin 类非零，即 $\int_M\Tr F\wedge F\ne0$。则对任意 $t\ne0,\infty$，方程 (3.29) 无解。最短的证明用\pA\anchor{src p.28, (3.32)}
\begin{equation}
\int_M\Tr F\wedge F
=\int_M\Tr\Bigl[\mathcal V_+(t)\wedge\mathcal V_+(-t^{-1})
+\mathcal V_-(t)\wedge\mathcal V_-(-t^{-1})\Bigr].
\label{eq:vt1}
\end{equation}
这是一个分部积分恒等式：右端是 $F$ 与 $D\varphi$ 的组合，展开后中间项全导数，剩下的正是 $\Tr F\wedge F$。于是若某个 $t$ 使 $\mathcal V_\pm(t)=0$，则右端为零，与 Pontryagin 类非零矛盾\pA。原文同时指出，这个证明在 $t=0,\infty$ 失效（$\mathcal V_\pm$ 有极点），而结论在那里也确实不成立：$t=0$ 与 $t=\infty$ 时方程至少有 $\varphi=0$ 的 instanton／anti-instanton 解\pA\anchor{src p.28}。

\textbf{定理 2（VT2）与恒等式 (3.33)。} 更强也更精细的恒等式是\pA\anchor{src p.29, (3.33)}
\begin{equation}
\begin{aligned}
&-\int_M \dd^4x\sqrt g\,\Tr\Bigl(\frac{t^{-1}}{t+t^{-1}}\mathcal V_+(t)_{\mu\nu}\mathcal V_+(t)^{\mu\nu}
+\frac{t}{t+t^{-1}}\mathcal V_-(t)_{\mu\nu}\mathcal V_-(t)^{\mu\nu}+\mathcal V_0^2\Bigr)\\
&\quad=-\int_M \dd^4x\sqrt g\,\Tr\Bigl(\tfrac12 F_{\mu\nu}F^{\mu\nu}
+D_\mu\varphi_\nu D^\mu\varphi^\nu+R_{\mu\nu}\varphi^\mu\varphi^\nu
+\tfrac12[\varphi_\mu,\varphi_\nu]^2\Bigr)
+\frac{t-t^{-1}}{t+t^{-1}}\int_M\Tr F\wedge F .
\end{aligned}
\label{eq:vt2}
\end{equation}
$R_{\mu\nu}$ 是 $M$ 的 Ricci 张量，来自分部积分把 $\int\sqrt g\,D_\mu\varphi_\nu D^\nu\varphi^\mu$ 与 $\int\sqrt g\,(D_\mu\varphi^\mu)^2$ 联系起来\pA。原文对它的评价是：
\begin{origquote}
What is surprising about this identity is that apart from the topological invariant proportional to $\int\Tr F\wedge F$, the right hand side is independent of $t$.
\end{origquote}
\noindent
\anchor{src p.29}。与 (3.32) 不同，(3.33) 在 $t\to0,\infty$ 有自然极限（把 $t(\mathcal V_-)^2/(t+t^{-1})$ 换成 $\lim_{t\to0}(t\mathcal V_-)^2$ 等），所以基于它的论证在 $t=0,\infty$ 也有效\pA\anchor{src p.29}。

\textbf{两个恒等式怎么用。} 我把它拆成三步（\pB，原文只用一句话带过）：
\begin{enumerate}[label=\arabic*)]
\item 若某 $t_0\ne0,\infty$ 有解，则由 (3.32) 得 $\int\Tr F\wedge F=0$。于是 (3.33) 的右端与 $t$ 无关。
\item 记 (3.33) 的左端为 $B(t)$。$t$ 实时系数 $t^{-1}/(t+t^{-1})=1/(1+t^2)$ 与 $t/(t+t^{-1})=t^2/(1+t^2)$ 都为正，所以 $B(t)$ 是平方和（带负号），即 $B(t)\ge0$。解存在给出 $B(t_0)=0$。
\item 由 (3.33) 与第 1 步，
\begin{equation}
B(u)-B(t)=\bigl(f(u)-f(t)\bigr)\int_M\Tr F\wedge F,
\qquad
f(t)=\frac{t-t^{-1}}{t+t^{-1}},
\label{eq:Bdiff}
\end{equation}
而 $\int\Tr F\wedge F=0$，故 $B(u)=B(t_0)=0$ 对一切 $u$ 成立；第 2 步的平方和说明每一项都为零，即 $\mathcal V_\pm(u)=\mathcal V_0=0$ 对一切 $u$ 成立。这就证明了\textbf{VT2}：某个 $t\ne0,\infty$ 的解自动是对所有 $t$ 的解，因而给出平坦 $G_{\mathbb C}$ 联络（$\mathcal V_\pm(t)=\mathcal V_0=0$ 对一切 $t$ 等价于 $\mathcal F=0$）\pA\anchor{src p.29}。
\end{enumerate}

\textbf{更强的 VT1$'$ 与 $f$ 的单调性。} 设 $\int\Tr F\wedge F>0$。若 $B(t)=0$，由 \eqref{eq:fmono} 得 $f(u)\ge f(t)$ 对一切 $u$ 成立，即 $f(t)$ 是 $f$ 的最小值；$\int\Tr F\wedge F<0$ 时同理是最大值。直接算
\begin{equation}
f'(t)=\frac{4t}{(1+t^2)^2},
\qquad
f(0)=-1,\quad f(\pm1)=0,\quad \lim_{t\to\pm\infty}f(t)=1,
\label{eq:fmono}
\end{equation}
（对实 $t$）可知 $f$ 在 $\mathbb R$ 上严格单调，唯一最小值在 $t=0$、唯一最大值在 $t\to\infty$。于是\pA\anchor{src p.30}
\begin{quote}
\textbf{VT1$'$}：$\int_M\Tr F\wedge F>0$ 时方程除 $t=0$ 外无解；$\int_M\Tr F\wedge F<0$ 时除 $t=\infty$ 外无解。
\end{quote}
这里的单调性计算原文只说「唯一最小／最大值」，我们把它写全了（\pB）。

\paragraph{结论的作用.}
vanishing theorems 把「一族 TFT」从修辞变成定理：对 $t\ne0,\infty$，解空间（因而无插入配分函数）与 $t$ 无关；Pontryagin 类非零的丛上干脆无解。这正是 M1 里「整个 $\mathbb{CP}^1$ 参数族保持拓扑性」的技术内容\pC\anchor{src p.28–30}。但要小心：本节\textbf{没有}证明与度规无关。度规无关性来自 §3.4 的结构性事实——作用量的度规变化是 $\delta_t$-恰当的\pB。

\caveat{(3.33) 在 $t=\pm\ii$ 处有极点（$t+t^{-1}=0$），所以整个 vanishing theorem 的论证只在实 $t$ 上有效；原文也明说「本节只考虑实 $t$」\pA\anchor{src p.28}。这不是疏漏：$t=\pm\ii$ 正是 B-模型式的点，那里的理论结构不同，要用另一套论证（§3.5 与 §5.2）。把 (3.33) 直接套到 $t=\ii$ 会得到 $0/0$ 型的假结论。}

\subsection{原文 §3.4：拓扑拉氏量——度规只出现在 $\delta_t$-恰当项里}

\begin{kcard}{技术路线卡：在弯曲 $M$ 上拼出拓扑作用量}
\anchor{src p.31–35}
\textbf{任务.} 找到作用量 $I$，使它对任意 $t$ 都被某个 $\delta_t$ 湮灭、且在 $M=\mathbb R^4$ 时回到 $\mathcal N=4$ SYM；由此得到度规无关性。

\textbf{输入.} 变换律 (3.23)–(3.28)、代数 $\delta_T^2=-\ii(u^2+v^2)\mathcal L_\sigma$、(3.33) 的恒等式。

\textbf{关键构造.} 三块拼装：$I_0=\delta_\ell\delta_r V$（只含 $A,\varphi,\sigma,\eta$ 的退化作用量）$+$ $I_1=\delta_t V_1$（补上规范动能与 $\chi$ 项）$+$ $I_2$（纯拓扑项，用来抵消 $I_1$ 残留的 $t$-依赖）。

\textbf{结论.} $I=I_0+I_1+I_2$ 与 $t$ 无关，因而对一切 $t$ 被 $\delta_t$ 湮灭；度规变化形如 $\delta_t(\cdots)$。

\textbf{出口.} 这一条是「拓扑性」的最终保证，S3 与 S5 中所有「与度规无关」的断言都以此为依据。

\textbf{锚点索引.} p.31 代数 (3.35)；p.32 辅助场 $P$ 与 $I_0$；p.33 $H$ 与 $V_1$；p.34 积掉 $H$ 得 (3.45)(3.46)；p.35 (3.49) 与总作用量；p.36 约化到 (2.10)。
\end{kcard}

\paragraph{动机.}
§3.3 的论证依赖一个前提：我们已经有一个在弯曲 $M$ 上定义的、对任意 $t$ 有拓扑对称的作用量。为什么不能直接把 $\mathbb R^4$ 上的 $\mathcal N=4$ 拉氏量「推广」过去？因为构造超对称作用量一般需要辅助场使代数离壳闭合，而 $\mathcal N=4$ 的辅助场问题是出了名的难（原文引 [53]，指出这是长期障碍）\pA\anchor{src p.32}。原文的策略因此不是「改造 $\mathcal N=4$ 拉氏量」，而是「用尽可能少的场、按 $\delta$-恰当的方式拼装」；拼装完成后，再验证它在 $\mathbb R^4$ 上回到 $\mathcal N=4$。

\paragraph{推导.}

\textbf{第一步：超代数。} 先算 $\delta_T^2$。对 $\mathcal K\ge0$ 的场\pA\anchor{src p.31, (3.35)}
\begin{equation}
\delta_T^2A=-\ii(u^2+v^2)(-D\sigma),\quad
\delta_T^2\varphi=-\ii(u^2+v^2)[\sigma,\varphi],\quad
\delta_T^2\psi=-\ii(u^2+v^2)[\sigma,\psi],
\label{eq:algebrasq}
\end{equation}
$\delta_T^2\tilde\psi$ 与 $\delta_T^2\sigma=0$ 同理。统一写法是
\begin{equation}
\delta_T^2\Phi=-\ii(u^2+v^2)\,\mathcal L_\sigma(\Phi),
\qquad
\mathcal L_\sigma(A)=-D\sigma,\quad \mathcal L_\sigma(\Phi)=[\sigma,\Phi]\ (\text{其余场}),
\label{eq:LT}
\end{equation}
即 $\delta_T^2$ 是 $\sigma$ 生成的规范变换乘以 $-u^2(1+t^2)=-(u^2+v^2)$。展开成 $u,v$ 的幂次等价于
\begin{equation}
\delta_\ell^2\Phi=\delta_r^2\Phi=-\ii\,\mathcal L_\sigma(\Phi),
\qquad
\{\delta_\ell,\delta_r\}\Phi=0 ,
\label{eq:lr}
\end{equation}
对 $\mathcal K<0$ 的场这个关系成立得要更曲折（部分场只在对运动方程成立时），但结论相同\pA\anchor{src p.31, (3.36)}。注意 \eqref{eq:LT} 在 $t=\pm\ii$ 时给出 $u^2+v^2=0$，即 $\delta_t^2=0$ 严格成立——这就是「$t=\pm\ii$ 退化」的第一层含义\pB。

\textbf{第二步：只对 $(\eta,\tilde\eta)$ 闭合的辅助场 $P$。} 引入李代数值零形式 $P$，并令\pA\anchor{src p.32, (3.37)}
\begin{equation}
\delta_T\sigma=\ii u\,\eta+\ii v\,\tilde\eta,\quad
\delta_T\eta=vP+u[\sigma,\bar\sigma],\quad
\delta_T\tilde\eta=-uP+v[\sigma,\bar\sigma],\quad
\delta_TP=-\ii v[\sigma,\eta]+\ii u[\sigma,\tilde\eta].
\end{equation}
$P$ 的引入只闭合了 $(\sigma,\eta,\tilde\eta,P)$ 这一小段代数，$\chi$ 仍不闭合——这是有意为之。接着取 $V$ 使 $P$ 的 Euler--Lagrange 方程给出 $P=D^*\varphi$，并令
\begin{equation}
I_0=\delta_\ell\delta_r V
=\frac{1}{e^2}\int_M\dd^4x\sqrt g\,\Tr\Bigl(\tfrac12 P^2-PD^*\varphi+\tfrac12[\sigma,\bar\sigma]^2
-D_\mu\bar\sigma D^\mu\sigma-[\varphi_\mu,\bar\sigma][\varphi^\mu,\sigma]+\cdots\Bigr),
\label{eq:I0}
\end{equation}
省略号是费米项，完整式见原文 (3.39)\pA\anchor{src p.32}。$I_0$ 有两个关键性质：（i）对任意 $t=v/u$ 都有拓扑对称，因为它只是 $\delta_\ell\delta_r V$ 而 $V$ 规范不变；（ii）度规\textbf{只}进入 $V$，不进入 $\delta_\ell,\delta_r$，所以度规变化是 $\delta_\ell\delta_r(\cdots)$ 形式——这正是 TFT 所需要的\pA\anchor{src p.33}。用 $P=D^*\varphi$ 可把 $\Tr(P^2-2PD^*\varphi)$ 换成 $-\Tr(D^*\varphi)^2$（(3.40)）。但 $I_0$ 是退化的：没有规范动能，也没有任何 $\chi$ 项\pA\anchor{src p.33}。

\textbf{第三步：为什么不能同时用 $\delta_\ell$ 和 $\delta_r$ 闭合 $\chi$。} 原文指出：可以为 $\chi$ 引入辅助场以闭合代数，但\textbf{无法}在同时保留 $\delta_\ell$ 与 $\delta_r$ 的前提下这样做；而 (3.33) 右端那个 $t$-依赖的拓扑项正是障碍\pA\anchor{src p.33}。于是改换策略：\textbf{固定一个 $t$}，只用 $\delta_t=\delta_\ell+t\delta_r$。引入二形式辅助场 $H$ 并规定\pA\anchor{src p.33, (3.41)–(3.42)}
\begin{equation}
\delta_t\chi=H,
\qquad
\delta_tH=-\ii(1+t^2)[\sigma,\chi],
\qquad
\delta_t^2\Phi=-\ii(1+t^2)\mathcal L_\sigma(\Phi)\quad(\Phi=\chi,H),
\label{eq:Halg}
\end{equation}
后一式正是 \eqref{eq:algebrasq} 在 $(u,v)=(1,t)$ 时的特化。为了让 $\chi$ 的变换与 §3.1 一致，运动方程必须给出
\begin{equation}
H^{+}=\mathcal V_+(t),
\qquad
H^{-}=t\,\mathcal V_-(t).
\label{eq:HEOM}
\end{equation}
取 $I_1=\delta_tV_1$，其中\pA\anchor{src p.33, (3.43)}
\begin{equation}
V_1=\frac{2}{e^2}\int_M\dd^4x\sqrt g\,\frac{1}{1+t^2}\Tr\Bigl(
\chi^{+}_{\mu\nu}\bigl(\tfrac12H^{+\mu\nu}-\mathcal V_+(t)^{\mu\nu}\bigr)
+\chi^{-}_{\mu\nu}\bigl(\tfrac12H^{-\mu\nu}-t\,\mathcal V_-(t)^{\mu\nu}\bigr)\Bigr).
\label{eq:V1}
\end{equation}
$\delta_tV_1$ 给出（原文 (3.44)）
\begin{equation}
I_1=\frac{1}{e^2}\int_M\dd^4x\sqrt g\,\Bigl[
\frac{2}{1+t^2}\Tr\bigl(\tfrac12H^+H^+-H^+\mathcal V_+(t)\bigr)
+\frac{2}{1+t^2}\Tr\bigl(\tfrac12H^-H^--tH^-\mathcal V_-(t)\bigr)
+(\text{费米项})\Bigr].
\label{eq:I1raw}
\end{equation}

\textbf{第四步：积掉 $H$（原文省略的全部步骤，\pB）。} 每个 $H$ 块都是二次式，配方即可：
\begin{equation}
\begin{aligned}
\frac{2}{1+t^2}\Bigl(\tfrac12 H^2-H\mathcal V_+\Bigr)
&=\frac{1}{1+t^2}(H-\mathcal V_+)^2-\frac{1}{1+t^2}\mathcal V_+^2,\\
\frac{2}{1+t^2}\Bigl(\tfrac12 H^2-tH\mathcal V_-\Bigr)
&=\frac{1}{1+t^2}(H-t\mathcal V_-)^2-\frac{t^2}{1+t^2}\mathcal V_-^2 .
\end{aligned}
\label{eq:completeSquare}
\end{equation}
极值点正是 \eqref{eq:HEOM} 的 $H^+=\mathcal V_+$、$H^-=t\mathcal V_-$；代回并使用恒等式
\begin{equation}
\frac{t^{-1}}{t+t^{-1}}=\frac{1}{1+t^2},
\qquad
\frac{t}{t+t^{-1}}=\frac{t^2}{1+t^2},
\label{eq:coeffid}
\end{equation}
就得到原文 (3.45) 的玻色部分 $-\frac{1}{e^2}\int\sqrt g\,\Tr\bigl(\frac{t^{-1}}{t+t^{-1}}\mathcal V_+^2+\frac{t}{t+t^{-1}}\mathcal V_-^2\bigr)$。这一步是纯代数，但它是整节唯一「原文说 integration 就跳过」的地方，值得写全。

\textbf{第五步：用 (3.33) 改写，残留一个拓扑项。} 把 (3.33) 代入 (3.45) 的玻色部分，得到（原文 (3.46)）\pA\anchor{src p.34}
\begin{equation}
\begin{split}
I_1={}&-\frac{1}{e^2}\int_M\dd^4x\sqrt g\,\Tr\Bigl(\tfrac12F_{\mu\nu}F^{\mu\nu}
+D_\mu\varphi_\nu D^\mu\varphi^\nu+R_{\mu\nu}\varphi^\mu\varphi^\nu
+\tfrac12[\varphi_\mu,\varphi_\nu]^2-(D^*\varphi)^2\Bigr)\\
&+(\text{费米项})
+\frac{t-t^{-1}}{e^2(t+t^{-1})}\int_M\Tr F\wedge F .
\end{split}
\label{eq:I1}
\end{equation}
除最后一项外，$I_1$ 与 $t$ 无关；最后一项是拓扑不变量，自动被 $\delta_t$（以及 $\delta_\ell,\delta_r$）湮灭，所以\textbf{可以再加一项来抵消它而不破坏拓扑对称}。

\textbf{第六步：加 $I_2$，选系数。} 原文取\pA\anchor{src p.35, (3.49)}
\begin{equation}
I_2=-\Bigl(\frac{t-t^{-1}}{e^2(t+t^{-1})}-\frac{\ii\theta}{8\pi^2}\Bigr)\int_M\Tr F\wedge F,
\label{eq:I2}
\end{equation}
于是 $I_1+I_2$ 的拓扑项只剩标准的 $\theta$ 项 $\frac{\ii\theta}{8\pi^2}\int\Tr F\wedge F$。总作用量
\begin{equation}
I=I_0+I_1+I_2
\label{eq:Itotal}
\end{equation}
的构造只保证它对\textbf{构造时用的那个} $t$ 被 $\delta_t$ 湮灭；但由于 $I$ 与 $t$ 无关，「对某个 $t$ 湮灭」自动升级为「对一切 $t$ 湮灭」\pA\anchor{src p.35}。

\paragraph{结论的作用.}
这一小节一次交出三件事，全是后面要反复引用的事实：
\begin{enumerate}[label=(\roman*)]
\item $I$ 有\textbf{完整的 $\mathbb{CP}^1$ 族拓扑对称}，因此闭四维流形上无插入的配分函数与 $t$ 无关\pA\anchor{src p.35}；
\item $I$ 是 TFT：$I_0,I_1$ 的度规依赖形如 $\delta_t(\cdots)$，$I_2$ 与度规无关，所以度规变化整体是 $\delta_t$-恰当的\pA\anchor{src p.35}；
\item 消去 $P$ 之后，玻色部分在 $M=\mathbb R^4$ 回到 §2 的 (2.10)，费米项由拓扑对称也对上，因此 $I$ 确实是 $\mathcal N=4$ SYM 到弯曲四维流形的推广\pA\anchor{src p.35}。
\end{enumerate}
第 (ii) 条是全文所有「拓扑不变量」断言的最终依据；第 (iii) 条保证我们没有偷偷换理论。

\caveat{$I_1$ 的分母 $1+t^2$ 与拓扑项的系数 $(t-t^{-1})/(t+t^{-1})$ 是两个不同的组合，别混：前者在 $t=\pm\ii$ 发散（$I_1$ 的构造在 $t=\pm\ii$ 失效），后者在 $t=\pm\ii$ 也有极点但可解析延拓到 $t=0,\infty$ 的极限。用 \eqref{eq:coeffid} 可以把前者写成 $1/(1+t^2)$、后者写成 $(t^2-1)/(t^2+1)$——这是检查符号最快的办法\pA\anchor{src p.34}。}

\subsection{原文 §3.5：canonical parameter $\Psi$——比 $\tau$ 更本质的参数}

\paragraph{动机.}
到这里，理论依赖两组参数：耦合 $\tau=\theta/2\pi+4\pi\ii/e^2$ 与扭化参数 $t$。原文明说它们\textbf{不独立}：
\begin{origquote}
But these parameters are not really all independent. We will now show that the theory really only depends on a certain combination $\Psi$ of $t$ and $\tau$, which we will call the canonical parameter.
\end{origquote}
\noindent\anchor{src p.35}。为什么这件事重要？因为它决定了 $S$-对偶怎么作用：如果 $\Psi$ 是真正内禀的参数，$S$-对偶就应该在 $\Psi$ 上作用为分式线性变换，而 $t$ 只是坐标。这也直接通向 §5 的二维解释（$\Psi$ 是复化 Kähler 类）。

\paragraph{推导.}

\textbf{第一步：把 $I$ 写成「$\delta_t$-恰当项 $+$ 拓扑项」。} 由 (3.49)，作用量形如\pA\anchor{src p.36, (3.51)}
\begin{equation}
I=\delta_tV+\frac{\ii\Psi}{4\pi}\int_M\Tr F\wedge F,
\qquad
-\frac{\Psi}{4\pi\ii}=\text{（拓扑项的系数）}.
\label{eq:Ipsi}
\end{equation}
只要 $\delta_tV$ 的细节无关紧要，可观测量的 $\tau,t$ 依赖就只通过系数 $\Psi$。但这里有一个表面的循环：$\delta_t$ 的定义本身依赖 $t$，所以「$I$ 只依赖 $\Psi$」还不能直接下结论。原文的处理（\pA\anchor{src p.36}）是：对 $t\ne\pm\ii$ 重新定义场，取 $\delta_t'=\delta_t/\sqrt{1+t^2}$，使
\begin{equation}
(\delta_t')^2=-\ii\,\mathcal L_\sigma
\label{eq:deltaprime}
\end{equation}
与 $t$ 无关；再把辅助场 $P,H$ 与费米子做小的重定义，每个多重态就变成标准形式 $\delta_t'U=V$、$\delta_t'V=-\ii\mathcal L_\sigma(U)$（cohomological field theory 的标准多重态，与等变上同调的 Cartan 模型有关）。于是变换律完全与 $t$ 无关，循环被打破\pA。

\textbf{第二步：从系数反解出 $\Psi$（完整推导）。} 由 \eqref{eq:I1}\eqref{eq:I2}，$I$ 中 $\int\Tr F\wedge F$ 的系数是
\begin{equation}
\frac{t-t^{-1}}{e^2(t+t^{-1})}+\frac{\ii\theta}{8\pi^2}=\frac{s}{e^2}+\frac{\ii\theta}{8\pi^2},
\qquad s=\frac{t-t^{-1}}{t+t^{-1}} .
\label{eq:coeff}
\end{equation}
用 $\mathrm{Re}\,\tau=\theta/2\pi$、$\mathrm{Im}\,\tau=4\pi/e^2$，有 $\ii\theta/(8\pi^2)=\ii\,\mathrm{Re}\,\tau/(4\pi)$、$s/e^2=s\,\mathrm{Im}\,\tau/(4\pi)$，于是
\begin{align}
\frac{s}{e^2}+\frac{\ii\theta}{8\pi^2}
&=\frac{\ii\,\mathrm{Re}\,\tau-s\,\mathrm{Im}\,\tau}{4\pi}\nonumber\\
&=\frac{\ii}{4\pi}\Bigl(\mathrm{Re}\,\tau+\ii s\,\mathrm{Im}\,\tau\Bigr)\nonumber\\
&=\frac{\ii\Psi}{4\pi}
=-\frac{\Psi}{4\pi\ii},
\label{eq:Psideriv}
\end{align}
其中最后一步要求
\begin{equation}
\boxed{\ \Psi=\mathrm{Re}\,\tau+\ii\,\mathrm{Im}\,\tau\,\frac{t-t^{-1}}{t+t^{-1}}
=\frac{\tau+\bar\tau}{2}+\frac{\tau-\bar\tau}{2}\cdot\frac{t-t^{-1}}{t+t^{-1}}\ }
\label{eq:Psi-twist}
\end{equation}
这正是原文 (3.50)\pA\anchor{src p.35}。反过来，把 \eqref{eq:Psi-twist} 通分可得一个更紧的形式（\pB，原文未写）
\begin{equation}
\Psi=\frac{\tau\,t^2+\bar\tau}{1+t^2},
\label{eq:Psit2}
\end{equation}
它把 $\Psi$ 显式写成 $t^2$ 的分式线性函数：给定 $\tau$，$t^2\mapsto\Psi$ 是 $\mathbb{CP}^1$ 的自同构。这个形式让下面的反解与 $S$-对偶推导都变成一行代数。

\textbf{第三步：反解 $t$（式 (3.52)）。} 由 \eqref{eq:Psit2}，
\begin{equation}
\Psi(1+t^2)=\tau t^2+\bar\tau
\quad\Longrightarrow\quad
t^2=\frac{\bar\tau-\Psi}{\Psi-\tau}=-\frac{\Psi-\bar\tau}{\Psi-\tau},
\label{eq:tinv}
\end{equation}
与原文 (3.52) 一致\pA\anchor{src p.36}。由 \eqref{eq:tinv} 可以立刻读出一批特殊值：$t=\pm\ii$（即 $t^2=-1$）时 $\Psi=\infty$，与 $\tau$ 无关；$t=\pm1$（即 $t^2=1$）时 $\Psi=\mathrm{Re}\,\tau$；$|t|=1$ 时 $t^2=\bar t^{-2}$ 使 $s$ 为实数，故 $\Psi$ 为实\pA\anchor{src p.36}。原文还指出：所有 $\Psi$ 都能在任意弱耦合（$\mathrm{Im}\,\tau$ 任意大）实现；固定 $\Psi$ 与 $\mathrm{Re}\,\tau$、令 $\mathrm{Im}\,\tau\to\infty$，由 \eqref{eq:tinv} 得 $t^2\to1$，即 $t\to\pm1$——弱耦合极限在参数空间里对应 $t=\pm1$ 的邻域\pA\anchor{src p.36–37}。

\textbf{第四步：$S$-对偶作用在 $\Psi$ 上（完整推导）。} 原文明说：把 $\tau$ 送到 $(a\tau+b)/(c\tau+d)$ 的变换也同样作用在 $\Psi$ 上\pA\anchor{src p.37, (3.53)}
\begin{equation}
\Psi\ \longmapsto\ \frac{a\Psi+b}{c\Psi+d}.
\label{eq:Spsi}
\end{equation}
检验 $T:\tau\to\tau+1$：$t$ 不变，故 $\Psi\to\Psi+1$ ✓。检验 $S:\tau\to-1/(n_{\mathfrak g}\tau)$：由 \eqref{eq:taction}，$t\to\pm t\tau/|\tau|$，故 $t^2\to t^2\tau/\bar\tau$。把这个与 $\tau'=-1/(n_{\mathfrak g}\tau)$ 代进 \eqref{eq:Psit2}：
\begin{align}
\Psi(\tau',t')
&=\frac{\tau't'^2+\bar\tau'}{1+t'^2}\nonumber\\
&=\frac{-\frac{1}{n_{\mathfrak g}\tau}\frac{t^2\tau}{\bar\tau}-\frac{1}{n_{\mathfrak g}\bar\tau}}{1+\frac{t^2\tau}{\bar\tau}}\nonumber\\
&=-\frac{1}{n_{\mathfrak g}}\cdot\frac{1+t^2}{\tau t^2+\bar\tau}\nonumber\\
&=-\frac{1}{n_{\mathfrak g}\Psi},
\label{eq:Sderiv}
\end{align}
最后一步用了 \eqref{eq:Psit2} 的 $\Psi$。这与原文 (3.54) 的四项式结果一致\pA\anchor{src p.37, (3.54)}；我们的写法（\pB）短得多，因为它把 (3.54) 里 $\tau,\bar\tau$ 的配对一次性吸收进 \eqref{eq:Psit2} 的通分形式。顺带解决了「$\Psi$ 会不会因 $t\to-t$ 而变」的问题：$\Psi$ 只依赖 $t^2$，所以 $t\to-t$ 在 $\Psi$ 上无作用——与「$t$ 与 $-t$ 等价」完全一致\pA\anchor{src p.37}。

\textbf{第五步：反定向对称。} 原文指出，$\mathrm{SL}(2,\mathbb Z)$ 作用在 $\mathbb{CP}^1$ 值的 $\Psi$ 上（不只是上半平面），所以可以把它扩充为 $\mathrm{GL}(2,\mathbb Z)$，加入\pA\anchor{src p.37, (3.55)}
\begin{equation}
\mathcal T=\begin{pmatrix}-1&0\\0&1\end{pmatrix},
\qquad
\mathcal T:\ \Psi\longmapsto-\Psi,
\label{eq:Trev}
\end{equation}
交换上下半平面。这个对称存在，当且仅当 $M$ 有一个反定向的对称：它保持 $e^2$（故 $\mathrm{Im}\,\tau$ 不变）、反转瞬子数（故 $\mathrm{Re}\,\tau$ 反号，即 $\tau\to-\bar\tau$），同时 $t\to1/t$\pA\anchor{src p.37}。用 \eqref{eq:Psit2} 直接验证（\pB）：
\begin{equation}
\Psi\bigl(-\bar\tau,\,1/t\bigr)
=\frac{-\bar\tau\,t^{-2}+\tau}{1+t^{-2}}
=\frac{-(\bar\tau+\tau t^2)}{1+t^2}=-\Psi(\tau,t).
\label{eq:Psiminus}
\end{equation}

图 \ref{fig:pspace} 把「$t$ 上的相因子旋转」与「$\Psi$ 上的分式线性变换」画在同一张参数空间图上，两者的差别一眼可见：$t$ 的轨道是绕原点的圆（不越过 $\pm\ii$），$\Psi$ 的轨道则被 $T$ 平移、被 $S$ 反演，并被反定向元素 $\mathcal T$ 翻到下半平面。

\begin{figure}[htbp]
\centering
\begin{tikzpicture}[font=\scriptsize,>=Stealth,scale=1.02]
% ================= (a) t 平面 =================
\begin{scope}
  \fill[blue!5] (-1.95,0) rectangle (1.95,1.70);
  \fill[orange!7] (-1.95,-1.70) rectangle (1.95,0);
  \draw[->] (-2.05,0)--(2.05,0) node[below] {$\mathrm{Re}\,t$};
  \draw[->] (0,-1.90)--(0,1.95) node[right] {$\mathrm{Im}\,t$};
  \draw[thick] (-1.95,0)--(1.95,0);
  \fill (0,0) circle (1.1pt); \node[below left] at (-0.03,-0.04) {$0$};
  \fill (1,0) circle (1.1pt); \node[below] at (1,-0.05) {$1$};
  \fill (-1,0) circle (1.1pt); \node[below] at (-1,-0.05) {$-1$};
  \node[below] at (1.86,-0.05) {$\infty$};
  \fill (0,1) circle (1.1pt); \node[right] at (0.06,1) {$\ii$};
  \fill (0,-1) circle (1.1pt); \node[right] at (0.06,-1) {$-\ii$};
  \draw[dashed,kwblue!75] (0,0) circle (0.55);
  \draw[dashed,kwblue!75] (0,0) circle (1.32);
  \node[kwblue!55!black] at (1.32,1.46) {$|t|>1$};
  \node[orange!60!black] at (1.32,-1.46) {$|t|<1$};
  \node[align=center] at (0,-2.22)
    {(a) $t\in\mathbb{CP}^1$：虚线为 $S$-对偶轨道 $|t|=\text{const}$};
\end{scope}
% ================= (b) Psi 平面 =================
\begin{scope}[xshift=6.0cm]
  \fill[green!5] (-1.65,0) rectangle (1.85,2.18);
  \fill[gray!10] (-1.65,-1.05) rectangle (1.85,0);
  \fill[kwblue!13] (0.5,0.866) arc (60:120:1) -- (-0.5,2.12) -- (0.5,2.12) -- cycle;
  \draw[thick] (0.5,0.866) arc (60:120:1);
  \draw[thick] (-0.5,0.866)--(-0.5,2.12);
  \draw[thick] (0.5,0.866)--(0.5,2.12);
  \draw[->] (-1.75,0)--(1.90,0) node[below] {$\mathrm{Re}\,\Psi$};
  \draw[->] (0,-1.15)--(0,2.22) node[right] {$\mathrm{Im}\,\Psi$};
  \node[kwblue!65!black] at (0,1.78) {$\mathcal D$};
  % T 的轨道
  \fill (0.15,1.05) circle (1.1pt); \node[above] at (0.15,1.12) {$\Psi_0$};
  \fill (1.15,1.05) circle (1.1pt); \node[below right,inner sep=1pt] at (1.16,1.00) {$\Psi_0\!+\!1$};
  \fill (-0.85,1.05) circle (1.1pt); \node[below left,inner sep=1pt] at (-0.86,1.00) {$\Psi_0\!-\!1$};
  \draw[->,dashed,kwblue!80] (0.27,1.13) to[bend left=14] (1.03,1.13);
  \draw[->,dashed,kwblue!80] (-0.03,1.13) to[bend right=14] (-0.73,1.13);
  % S 反演
  \fill (-0.13,0.93) circle (1.1pt);
  \draw[->,dashed,kwred!85] (0.06,0.94) to[bend right=45] node[below right,inner sep=1pt] {$S$} (-0.10,0.84);
  % 反定向（上下半平面互换）
  \draw[->,dashed,kwred!70] (-1.46,1.82) to[bend left=24] node[left,inner sep=1pt] {$\mathcal T$} (-1.46,-0.82);
  \node[align=center] at (0,-1.35)
    {(b) $\Psi\in\mathbb{CP}^1$：阴影为 $\mathrm{SL}(2,\mathbb Z)$ 基本域 $\mathcal D$};
\end{scope}
\end{tikzpicture}
\caption{参数空间上的 $S$-对偶作用。(a) $t$ 平面上 $S$-对偶把 $t$ 乘一个相因子（式 \eqref{eq:taction}），轨道是 $|t|=$ 常数的圆；$\pm\ii$ 与实轴分别是 B-模型点与 A-模型区域。(b) $\Psi$ 平面上，$T:\Psi\mapsto\Psi+1$ 把 $\Psi_0$ 平移到 $\Psi_0\pm1$，$S:\Psi\mapsto-\frac{1}{n_{\mathfrak g}\Psi}$ 把 $\Psi_0$ 送到单位圆附近（式 \eqref{eq:Spsi}），反定向元素 $\mathcal T:\Psi\mapsto-\Psi$ 交换上下半平面（式 \eqref{eq:Trev}）。原文未画此图，按 \eqref{eq:taction}\eqref{eq:Spsi}\eqref{eq:Trev} 绘制\pB\anchor{src p.22–23, p.37}。}
\label{fig:pspace}
\end{figure}

\paragraph{结论的作用.}
$\Psi$ 是本节真正的「物理参数」：$I$ 只依赖 $\Psi$（在 $\delta_tV$ 无关紧要的意义上）；$S$-对偶在 $\Psi$ 上作用为全纯的分式线性变换 \eqref{eq:Spsi}；$T$ 是平移、$S$ 是 $-\frac{1}{n_{\mathfrak g}\Psi}$ 型的反演；反定向对称给出 $\Psi\to-\Psi$。这就是主链上「$S$-对偶为什么能作用在扭化族上」的全部根据。$t=\pm\ii$ 是 $\Psi=\infty$ 的退化点：那里 $\Psi$ 的唯一参数论证因极点失效，理论变成复结构 $\pm J$ 的 B-模型，$\tau$ 完全不起作用（原文说 Marcus [53] 直接证明了这一点）\pA\anchor{src p.36}。

\begin{warn}
\textbf{排除错误直觉（六）：$\Psi$ 就是 $\frac{t-\ii}{t+\ii}$ 之类的随手代换。} 错，而且错得很结构性。$(t-\ii)/(t+\ii)$ 是一个只依赖 $t$ 的交叉比：它把 $\pm\ii$ 送到 $0,\infty$，除此之外与 $\tau$ 毫无关系。但 $\Psi$ \textbf{按定义是作用量里拓扑项的系数}（\eqref{eq:Ipsi}），它必然含 $\tau$：由 \eqref{eq:Psi-twist}，$\Psi=\mathrm{Re}\,\tau+\ii\,\mathrm{Im}\,\tau\,s$。两者只在两个性质上巧合相同：$|t|=1$ 时 $\Psi$ 实、$t=\pm\ii$ 时 $\Psi=\infty$。差别在关键处：$S$-对偶作用在 $\Psi$ 上是\textbf{全纯}分式线性变换 \eqref{eq:Spsi}，而作用在 $t$ 上带 $|\cdot|$（式 \eqref{eq:taction}），不是全纯的。把 $\Psi$ 当成 $t$ 的函数而不含 $\tau$，就丢掉了整个 $S$-对偶结构。原文从未给出这样的公式；它是读者最容易「顺手猜」错的地方\pC。
\end{warn}

\begin{table}[htbp]
\centering\footnotesize
\caption{$t$ 的特殊值与四维／二维含义的对照。四维列为原文 §3.2 的方程性质；二维列为 §5.1 的对应（$w_+=-t$，$w_-=t^{-1}$）；$\Psi$ 列由 \eqref{eq:Psi-twist} 读出\pA\anchor{src p.26–28, p.36, p.61–62}。}
\label{tab:special}
\begin{tabular}{@{}lccc l@{}}
\toprule
\tabhead{$t$} & \tabhead{四维：方程 (3.29)} & \tabhead{二维模型} & \tabhead{$\Psi$} & \tabhead{备注}\\
\midrule
$0$ & $(D\varphi)^{-}=0$，含 instanton 解 & A-模型（复结构 $I$ 或 $-I$） & $\bar\tau$ & 实椭圆族的端点\\
$\pm1$ & 实椭圆，最「对称」 & A-模型（复结构 $\pm K$） & $\mathrm{Re}\,\tau$ & 弱耦合极限 $t\to\pm1$\\
实 $t$ & 实椭圆方程族 & A-模型 & $\mathrm{Re}\tau+\ii\mathrm{Im}\tau\,s$ & $s\in(-1,1)$\\
$|t|=1$ & 同上（$t=e^{\ii\vartheta}$） & A-模型 & 实 & 只靠 $\theta$ 调 $\Psi$\\
$\pm\ii$ & $\mathcal F=0$：平坦 $G_{\mathbb C}$ 联络 & B-模型（复结构 $\pm J$） & $\infty$ & $\tau$ 完全无关\\
$\infty$ & $(D\varphi)^{+}=0$，含 anti-instanton 解 & A-模型（复结构 $-I$ 或 $I$） & $\tau$ & 实椭圆族的另一端点\\
\bottomrule
\end{tabular}
\end{table}

\subsection{原文 §5.1：二维化——扭化 TFT 与 $\mathbb{CP}^1_g\hookrightarrow\mathbb{CP}^1_h\times\mathbb{CP}^1_h$}

\paragraph{动机.}
四维的族 $\mathbb{CP}^1_g$ 是在任意四维流形上定义的。到了二维，为什么还要重新讲一遍？因为（i）$S$-对偶的几何解释（镜像对称）只在二维才看得见；（ii）二维的构造会给出一个\textbf{更大}的族 $\mathbb{CP}^1_h\times\mathbb{CP}^1_h$，而四维只给其中一个一维子族。搞清楚「四维族嵌在哪里」，是 S4 与 S5 的技术前提。

\paragraph{推导.}

\textbf{第一步：二维扭化 TFT 的构造。} 从目标流形 $X$ 的 sigma 模型出发，选一对复结构 $(J_+,J_-)$ 使模型有 $(2,2)$ 超对称（原文引 [66] 的条件；该条件后来被 [67] 用广义复几何重新解释）；有 $(2,2)$ 之后有标准扭化配方 [70]\pA\anchor{src p.58}。对 $\MH$ 而言 $X$ 是超凯勒的，因而有一整族以 $w\in\mathbb{CP}^1_h$ 参数化的复结构 $I_w$；取 $(J_+,J_-)=(I_{w_+},I_{w_-})$ 自动满足 $(2,2)$ 条件\pA\anchor{src p.58}。扭化后的 TFT 由未破缺超对称方程刻画：玻色场 $\Phi:\Sigma\to X$，
\begin{equation}
(1-\ii J_+)\partial\Phi=0,
\qquad
(1-\ii J_-)\bar\partial\Phi=0,
\label{eq:twodim}
\end{equation}
它们分别来自费米场 $\chi_+,\chi_-$ 的 $\{Q,\chi_+\}=(1-\ii J_+)\partial\Phi$、$\{Q,\chi_-\}=(1-\ii J_-)\bar\partial\Phi$；第一式说 $\Phi$ 在 $J_+$ 下全纯，第二式说在 $J_-$ 下反全纯\pA\anchor{src p.59}。原文点出一个反直觉的事实：在紧致 Euclidean $\Sigma$ 上，$J_-\ne J_+$ 时 generic 的解是常映射，解与 $(J_+,J_-)$ 无关，但\textbf{方程}依赖 $(J_+,J_-)$——信息在方程里，不在解里\pA\anchor{src p.59}。

\textbf{第二步：A／B 模型是特例。} 若 $J_+=J_-=\hat J$，(5.5)(5.6) 合并为 $(1-\ii\hat J)\dd\Phi=0$，即 $\dd\Phi=0$：常映射，这就是复结构 $\hat J$ 的 B-模型。若 $J_+=-\bar J_-=\hat J$，就是 A-模型（实际依赖与 $\hat J$ 相配的辛形式 $\omega_{\hat J}$，而非复结构本身）；在 Euclidean $\Sigma$ 上 A-模型的两个方程冗余（互为复共轭），任一都断言 $\Phi$ 在 $\hat J$ 下全纯\pA\anchor{src p.59}。若 $J_-\ne\pm J_+$，则既非 A 也非 B——这正是需要广义复几何的地方（§5.2）。

\textbf{第三步：$\mathbb{CP}^1_h$ 上的复结构显式公式。} 超凯勒流形上的复结构族是\pA\anchor{src p.60, (5.7)}
\begin{equation}
I_w=\frac{1-w\bar w}{1+w\bar w}\,I+\frac{\ii(w-\bar w)}{1+w\bar w}\,J+\frac{w+\bar w}{1+w\bar w}\,K,
\qquad w\in\mathbb{CP}^1_h,
\label{eq:Iw-twist}
\end{equation}
其中 $I,J,K$ 是超凯勒三元组；$w=\pm\ii$ 时 $I_w=\mp J$。$I_w$ 的全纯坐标是 $A_z-w\varphi_z$ 与 $A_z+w^{-1}\varphi_z$\pA\anchor{src p.60}。

\textbf{第四步：从四维降到二维——用 abelian 情形做探针。} 四维族 $\mathbb{CP}^1_g$ 要求 $Q$ 在完整的 $\mathrm{Spin}'(4)$ 下不变（对任意 $M$ 有效）；当 $M=\Sigma\times C$ 时 holonomy 降到 $\mathrm{SO}(2)\times\mathrm{SO}(2)$，只需 $Q$ 在这些更小的群下不变，于是族变大。原文先在 abelian 情形（$G$ 破到极大环面）里确定 $\mathbb{CP}^1_g$ 的嵌入函数 $w_\pm(t)$\pA\anchor{src p.60–61}。abelian 方程是\pA\anchor{src p.60, (5.8)}
\begin{equation}
(F+t\,\dd\varphi)^{+}=0,
\qquad
(F-t^{-1}\dd\varphi)^{-}=0,
\label{eq:abelian}
\end{equation}
其中 $\varphi$ 是 $T^*M=T^*\Sigma\oplus T^*C$ 的截面。取 $\Sigma$ 的复坐标 $y$、$C$ 的复坐标 $z$，用 $\dd y\wedge\dd z$ 自对偶、$\dd y\wedge\dd\bar z$ 反自对偶，两式分别给出（原文 (5.9)(5.10)）\pA\anchor{src p.61}
\begin{equation}
\partial_y\bigl(A_z+t\,\varphi_z\bigr)=0,
\qquad
\bar\partial_y\bigl(A_{\bar z}-t^{-1}\varphi_{\bar z}\bigr)=0 .
\label{eq:reduced}
\end{equation}
与 (5.5)(5.6) 比较：若 $A_z+t\varphi_z$ 在某个复结构 $I_{w_+}$ 下全纯，则 $w_+=-t$；若 $A_{\bar z}-t^{-1}\varphi_{\bar z}$ 在 $I_{w_-}$ 下全纯，则 $w_-=t^{-1}$。故\pA\anchor{src p.61, (5.11)}
\begin{equation}
w_+=-t,
\qquad
w_-=t^{-1}.
\label{eq:wpm}
\end{equation}
两个立刻可读的推论：B-模型要求 $w_+=w_-$，即 $t^{-1}=-t$，故 $t=\pm\ii$（此时 $I_w$ 为 $J$ 或 $-J$）；A-模型要求 $I_{w_+}$ 与 $I_{w_-}$ 相反，而 $w\mapsto-1/w$ 把 $I_w$ 送到其相反，故 $t=\bar t$，即实 $t$ 给 A-模型\pA\anchor{src p.61}。

\textbf{第五步：更大的族与唯一不变量。} $M=\Sigma\times C$ 时两个 $\mathbb{CP}^1_h$ 各自独立，得到 $\mathbb{CP}^1_h\times\mathbb{CP}^1_h$。流形的微分同胚群给出 $\mathbb{C}^*$ 作用 $w\to\lambda^{-1}w$，$(w_+,w_-)$ 与 $(\lambda^{-1}w_+,\lambda^{-1}w_-)$ 给同一个 TFT，唯一不变量是\pA\anchor{src p.62}
\begin{equation}
q=\frac{w_+}{w_-}=-t^2 .
\label{eq:q}
\end{equation}
$t$ 与 $-t$ 给二维等价的理论，正对应 (3.14)。注意不是所有 $\mathbb{C}^*$ 轨道都有四维代表：不动点 $(0,0)$、$(\infty,\infty)$ 与只有一个坐标为 $0/\infty$ 的轨道不在 $\mathbb{CP}^1_g$ 上；其中 $(0,0)$、$(\infty,\infty)$ 是复结构 $I$ 与 $-I$ 的 B-模型\pA\anchor{src p.62}。原文强调：复结构 $I$ 的 B-模型是数学上解释几何 Langlands 的起点（Hitchin 纤维化在 $I$ 下全纯，纤维上的 $T$-对偶与四维 $S$-对偶相关），但它\textbf{不}在四维族里；本文依赖的是四维的 $\mathbb{CP}^1_g$\pA\anchor{src p.62}。

\textbf{第六步：$S$-对偶作用在 $(w_+,w_-)$ 上。} 由 \eqref{eq:taction}，$t\to t(c\tau+d)/|c\tau+d|$；$t',t''$ 同样变换，结合 (5.20) 得\pA\anchor{src p.64, (5.22)}
\begin{equation}
w_+\longmapsto w_+\,\frac{c\tau+d}{|c\tau+d|},
\qquad
w_-\longmapsto w_-\,\frac{|c\tau+d|}{c\tau+d}.
\label{eq:Sww}
\end{equation}
重要特例：$B$-模型 $(0,0)$ 与 $A$-模型 $(0,\infty)$ 在 $S$-对偶下完全不变\pA\anchor{src p.64}。

\paragraph{结论的作用.}
$\mathbb{CP}^1_g$ 被具体钉在 $\mathbb{CP}^1_h\times\mathbb{CP}^1_h$ 里：曲线 $(w_+,w_-)=(-t,t^{-1})$，即 $q=-t^2$。这条曲线上的点才有四维解释；曲线外的点（如 $(0,0)$）是二维的「辅助」TFT，虽然不在四维族里，却在 S4 里扮演镜像对称的端点。原文据此说明：三维 $\to$ 二维的约化给出了一个更大族，而几何 Langlands 是关于其中\textbf{特定一维子族}的陈述\pA\anchor{src p.64}。

\begin{warn}
\textbf{排除错误直觉（七）：四维的 $\mathbb{CP}^1$ 就是二维的全部参数空间。} 错。二维有\textbf{两个}参数 $(w_+,w_-)$，四维只给出一维子族 $q=w_+/w_-=-t^2$；多出来的方向来自「只需要 $\mathrm{Spin}'(2)\times\mathrm{Spin}'(2)$ 不变」这个更弱的条件（§5.1 的 (5.13)–(5.17)）\pA\anchor{src p.63}。把两者混同，就会误以为 $(0,0)$ 处的 B-模型也来自四维——恰恰相反，原文明确指出 $(0,0)$ 不在 $\mathbb{CP}^1_g$ 上\pA\anchor{src p.62}。
\end{warn}

\caveat{(5.11) 的符号与指数不要抄反：$w_+=-t$，$w_-=t^{-1}$（不是 $w_+=t$、$w_-=-t^{-1}$）。检验办法：$t=\ii$ 时应有 $w_+=w_-$（B-模型），而 $-\ii=\ii^{-1}$ ✓；$t=-\ii$ 时 $w_+=\ii$、$w_-=-\ii$，对应 $I_{\pm\ii}=\mp J$，也对。任何别的符号约定都通不过这两个检验。}

\subsection{原文 §5.2：广义复几何——A／B 两型的统一语言}

\paragraph{动机.}
在二维，A-模型依赖辛形式、B-模型依赖复结构，看上去是两类不同的 TQFT。但四维化约给出的族在 $t$ 从实变到 $\pm\ii$ 时「从 A 变 B」，说明两者应是同一个几何对象的两端。Hitchin 的广义复几何正是这个对象：它在 $T\oplus T^*$ 上定义「复结构」，把辛结构与复结构作为两个特例，并用 $B$-场变换连接它们\pA\anchor{src p.65}。此外它还能回答一个关键问题：为什么四维化约出来的族不依赖 $C$ 的复结构？

\paragraph{推导.}

\textbf{第一步：定义。} 令 $\hat T=TX\oplus T^*X$，截面记作 $(v,\xi)$，$v\in TX$、$\xi\in T^*X$。$\hat T$ 上有自然的不定度规
\begin{equation}
(v,\xi)\cdot(v',\xi')=v^i\xi'_i+v'^i\xi_i,
\label{eq:gcMetric}
\end{equation}
使 $TX$ 与 $T^*X$ 都是迷向子空间。\term{generalized complex structure}{广义复结构}是一个线性变换 $\mathcal I:\hat T\to\hat T$，满足 $\mathcal I^2=-1$ 且保度规 \eqref{eq:gcMetric}，并满足适当的可积性条件（等价说法：$\mathcal I$ 的 $+\ii$ 本征分布对 Courant 括号封闭）\pA\anchor{src p.65}。两个基本例子：
\begin{equation}
\mathcal I_{\hat J}=\begin{pmatrix}\hat J&0\\0&-\hat J^{t}\end{pmatrix}
\quad(\text{B-模型}),
\qquad
\mathcal I_\omega=\begin{pmatrix}0&-\omega^{-1}\\ \omega&0\end{pmatrix}
\quad(\text{A-模型}),
\label{eq:gcExamples}
\end{equation}
分别由复结构 $\hat J$ 与辛形式 $\omega$ 给出\pA\anchor{src p.66, (5.23)–(5.24)}。$B$-场的进入方式是：对闭二形式 $B_0$，令
\begin{equation}
\mathcal M(B_0)=\begin{pmatrix}1&0\\ B_0&1\end{pmatrix},
\qquad
\mathcal I(B_0)=\mathcal M(B_0)\,\mathcal I\,\mathcal M(B_0)^{-1},
\label{eq:Bfield}
\end{equation}
后者仍可积，于是得到「带 $B$-场的 B-模型／A-模型」\pA\anchor{src p.66, (5.25)–(5.26)}。原文还转述 [67] 第 6 章的结果：sigma 模型的 $(2,2)$ 超对称等价于存在一对相容的广义复结构；每个扭化 TFT 由其中一个广义复结构决定，与第二个无关（Kähler 且 $B$ 平坦时，两个分别给出 B-模型与 A-模型；有 [72,68] 的支持，但没有完整证明）\pA\anchor{src p.66}。

\textbf{第二步：把四维族翻译成 $\mathcal I_t$。} 我们的情形是超凯勒 $X=\MH$，$B$-场先取零（之后可用 $B$-场变换恢复），$\omega_\pm=gJ_\pm$，$(J_+,J_-)=(I_{w_+},I_{w_-})$。由 [67] 的 eq.\ 6.3，两个相容广义复结构之一是\pA\anchor{src p.67, (5.27)}
\begin{equation}
\mathcal J=\frac12\begin{pmatrix}J_++J_-&-(\omega_+^{-1}-\omega_-^{-1})\\ \omega_+-\omega_-&-(J_+^{t}+J_-^{t})\end{pmatrix}.
\label{eq:Jgc}
\end{equation}
代入 (5.7) 与 $(w_+,w_-)=(-t,t^{-1})$，得\pA\anchor{src p.67, (5.28)–(5.29)}
\begin{equation}
\begin{aligned}
J_+&=\frac{1-t\bar t}{1+t\bar t}\,I-\frac{\ii(t-\bar t)}{1+t\bar t}\,J-\frac{t+\bar t}{1+t\bar t}\,K,\\
J_-&=-\frac{1-t\bar t}{1+t\bar t}\,I-\frac{\ii(t-\bar t)}{1+t\bar t}\,J+\frac{t+\bar t}{1+t\bar t}\,K,\\
\omega_+&=\mathrm{Im}\,\tau\Bigl(\frac{1-t\bar t}{1+t\bar t}\omega_I-\frac{\ii(t-\bar t)}{1+t\bar t}\omega_J-\frac{t+\bar t}{1+t\bar t}\omega_K\Bigr),\\
\omega_-&=\mathrm{Im}\,\tau\Bigl(-\frac{1-t\bar t}{1+t\bar t}\omega_I-\frac{\ii(t-\bar t)}{1+t\bar t}\omega_J+\frac{t+\bar t}{1+t\bar t}\omega_K\Bigr),
\end{aligned}
\label{eq:Jpm}
\end{equation}
其中 $\mathrm{Im}\,\tau$ 因子来自原文 (4.15)（超凯勒度规的归一化）。代回 (5.27) 并利用「超凯勒上 $a^2+b^2+c^2=1$ 时 $(a\omega_I+b\omega_J+c\omega_K)^{-1}=a\omega_I^{-1}+b\omega_J^{-1}+c\omega_K^{-1}$」，得到\pA\anchor{src p.67, (5.30)}
\begin{equation}
\mathcal I_t=\frac{1}{1+t\bar t}
\begin{pmatrix}
-\ii(t-\bar t)J & -(\mathrm{Im}\,\tau)^{-1}\bigl[(1-t\bar t)\omega_I^{-1}-(t+\bar t)\omega_K^{-1}\bigr]\\[4pt]
\mathrm{Im}\,\tau\bigl[(1-t\bar t)\omega_I-(t+\bar t)\omega_K\bigr] & \ii(t-\bar t)J^{t}
\end{pmatrix}.
\label{eq:It}
\end{equation}
\textbf{第一个重要观察：$I$、$K$、$\omega_J$ 全部消失。} $\mathcal I_t$ 只依赖 $J$、$\omega_K$、$\omega_I$，等价地只依赖 $J$ 与全纯二形式 $\Omega_J=\omega_K+\ii\omega_I$。这组数据在 $\MH$ 上与 $C$ 的度规无关（§4.1 的结论）——于是\textbf{从二维的角度解释了四维化约族为什么不依赖 $C$ 的复结构}\pA\anchor{src p.67}。这正是原文 ``Dependence On The Metric'' 一段想说的话。

\textbf{第三步：generic 的 $\mathcal I_t$ 是辛结构经 $B$-场变换。} 原文引 [67] §4.6 的 Prop.\ 4.34：这类族里的 generic 元素是某个辛结构导出的广义复结构的 $B$-场变换。具体地，存在 $\omega_0,B_0$ 使\pA\anchor{src p.68, (5.31)–(5.32)}
\begin{equation}
\mathcal I_t=\mathcal M(B_0)\begin{pmatrix}0&-\omega_0^{-1}\\ \omega_0&0\end{pmatrix}\mathcal M(B_0)^{-1},
\label{eq:Isym}
\end{equation}
\begin{equation}
\begin{aligned}
\omega_0&=\mathrm{Im}\,\tau\cdot\frac{1-\bar t^2t^2}{(1+t^2)(1+\bar t^2)}\Bigl(\omega_I-\frac{t+\bar t}{1-t\bar t}\,\omega_K\Bigr),\\
B_0&=-\mathrm{Im}\,\tau\cdot\frac{\ii(t^2-\bar t^2)}{(1+t^2)(1+\bar t^2)}\Bigl(\omega_I+\frac{1-t\bar t}{t+\bar t}\,\omega_K\Bigr).
\end{aligned}
\label{eq:omega0}
\end{equation}
因此对 $t\ne\pm\ii$，$\mathcal I_t$ 定义的 TQFT 等价于\textbf{辛形式 $\omega_0$、$B$-场 $B_0$ 的 A-模型}。这就在二维里解释了 canonical parameter $\Psi$：A-模型只依赖上同调类 $[B_0+\ii\omega_0]$（Kähler 上叫作复化 Kähler 类）。本问题里 $\omega_K$ 的上同调类为零，故\pA\anchor{src p.68, (5.33)}
\begin{equation}
[B_0+\ii\omega_0]=-\ii\,\mathrm{Im}\,\tau\,\frac{t-t^{-1}}{t+t^{-1}}\,[\omega_I].
\label{eq:class}
\end{equation}
\textbf{第四步：恢复四维 $\theta$ 角，得到 $\Psi$ 的二维意义。} 上面的 $\mathcal I_t$ 是从 $\theta=0$ 出发的；四维 $\theta$ 角给额外贡献 $B'=-(\theta/2\pi)\omega_I=-(\mathrm{Re}\,\tau)\omega_I$（§4.1 末），用 $\mathcal M(B')$ 共轭并入即可。于是模型依赖\pA\anchor{src p.68}
\begin{equation}
[B_0+B'+\ii\omega_0]
=-[\omega_I]\Bigl(\mathrm{Re}\,\tau+\ii\,\mathrm{Im}\,\tau\,\frac{t-t^{-1}}{t+t^{-1}}\Bigr)
=-[\omega_I]\,\Psi .
\label{eq:Psired}
\end{equation}
这就是 (3.50) 的二维解释：\textbf{$\Psi$ 是四维扭化族在二维里的复化 Kähler 类}；$S$-对偶作用在 $\Psi$ 上的分式线性变换，就是二维镜像对称的通常作用\pA\anchor{src p.68, (5.34)}。

\paragraph{结论的作用.}
广义复几何把本节的所有线索收在一处：（i）A／B 模型不是两类理论，而是同一个 $\mathcal I_t$ 族的两个极限；（ii）族的一维子族 $\mathbb{CP}^1_g$ 由 (5.30) 显式给出，参数 $t$ 就是 $B$-场变换的参数；（iii）$\Psi$ 有了二维身份——复化 Kähler 类，于是 $S$-对偶的几何意义（镜像对称）在二维浮现；（iv）$C$-度规无关性变成「$\mathcal I_t$ 只依赖 $J$ 与 $\Omega_J$」这一句话\pA\anchor{src p.67–68}。M2（四维 $S$-对偶在二维变成镜像对称）的技术内容在此成型\pC。

\caveat{(5.32) 里 $\omega_0$ 与 $B_0$ 的系数是两种不同的组合：$\omega_0$ 带 $(1-\bar t^2t^2)$、$B_0$ 带 $\ii(t^2-\bar t^2)$，且 $B_0$ 的第二项分母是 $t+\bar t$（在 $t=\pm\ii$ 发散）。原文没有给推导，只引 [67] 的命题；抄写时最容易错的是 $B_0$ 的 $\ii$ 因子与 $t+\bar t$ 的分母。检验办法：$t$ 实时 $t^2-\bar t^2=0$，故 $B_0=0$，$\omega_0$ 退化为实 $t$ 的 A-模型辛形式——这与 §5.1 的实 $t\leftrightarrow$ A-模型一致\pB\anchor{src p.68}。}

\begin{warn}
\textbf{补充：原文对广义复几何路线的态度。} 原文明确承认：对实 $t$，TQFT 本来就是 A-模型，上述论证并不真的\textbf{需要}广义复几何；作者只是相信它是自然的框架，并承认当前理解仍有技术缺口\pA\anchor{src p.69}。我们判断（\pC）：它的价值在于给出\textbf{统一的几何语言}——它让 $\mathbb{CP}^1_g$ 的度规无关性、$\Psi$ 的二维意义、$S$-对偶的镜像对称解释成为同一件事的三面，这对 S4 与 S5 是必需的。
\end{warn}

\subsection{小结：主链 S2 的记账与自检}

把本节的技术动作与它们的消费者对齐：

\begin{table}[htbp]
\centering\footnotesize
\caption{主链 S2 的技术动作—结论—消费者对照。右列为下游使用者（S3 = 原文 §4 紧化；S4 = §5.3–5.6；S5 = §6 线算子）\pC}。
\label{tab:ledger}
\begin{tabular}{@{}p{3.9cm}p{6.2cm}p{4.3cm}@{}}
\toprule
\tabhead{技术动作} & \tabhead{结论} & \tabhead{消费者}\\
\midrule
超荷全是旋量 & 必须有标量超荷才能谈拓扑 & 整节的前提\\
对角嵌入 $\mathrm{Spin}'(4)$ & 恰好两个不变超荷 $\varepsilon_\ell,\varepsilon_r$ & $\mathbb{CP}^1$ 族\\
$\mathcal K$ 荷相同 & $t$ 是真参数；$t\sim-t$ 是唯一平凡等价 & $S$-对偶在 $t$ 上的作用\\
固定点方程 (3.29) & $t=\pm\ii$：平坦 $G_{\mathbb C}$ 联络；实 $t$：椭圆族 & S3 的 Hitchin 方程\\
vanishing theorems & 解空间（$t\ne0,\infty$）不随 $t$ 变 & M1\\
拓扑拉氏量 & 度规变化 $=\delta_t(\cdots)$；$\mathbb R^4$ 上回到 $\mathcal N=4$ & 全部拓扑不变量\\
canonical parameter & 只依赖 $\Psi$；$S$-对偶 $=$ 分式线性变换 & S4 的镜像对称\\
$\mathbb{CP}^1_g\subset\mathbb{CP}^1_h\times\mathbb{CP}^1_h$ & $(w_+,w_-)=(-t,t^{-1})$，$q=-t^2$ & S4 的二维族\\
广义复几何 & $\mathcal I_t$ 只依赖 $J,\Omega_J$；$\Psi=[B+\ii\omega]$ & S4、S5 的统一语言\\
\bottomrule
\end{tabular}
\end{table}

\noindent\textbf{自检问题.}
\begin{enumerate}[label=\arabic*.,itemsep=2pt,topsep=3pt]
\item 为什么「两个不变超荷的 $\mathcal K$ 相同」是 $t$ 能成为物理参数的关键？若差一个 $\mathcal K$，$S$-对偶在 $t$ 上会变成什么作用？§5 的镜像对称还成立吗？
\item 用 \eqref{eq:Psi-twist} 的通分形式 \eqref{eq:Psit2} 重新推导：$t=\pm1$ 时 $\Psi=\mathrm{Re}\,\tau$；$t=\pm\ii$ 时 $\Psi=\infty$；并说明为什么 $t\to-t$ 在 $\Psi$ 上无作用，而 $t\to1/t$ 结合 $\tau\to-\bar\tau$ 给出 $\Psi\to-\Psi$。
\item (3.33) 为什么在 $t=0,\infty$ 有自然极限、在 $t=\pm\ii$ 没有？把两个极限的系数分别写出来，并说明它们对应的物理极限（instanton／anti-instanton 解）。
\item 在 $\mathbb{CP}^1_h\times\mathbb{CP}^1_h$ 中，哪些轨道不来自四维族？$(0,0)$ 处的 TFT 是什么？为什么它不在 $\mathbb{CP}^1_g$ 上，却仍对几何 Langlands 重要？
\end{enumerate}

\enlargethispage{\baselineskip}
\noindent\textbf{本章一句话.} 扭化不是把理论变简单，而是\textbf{换对角嵌入、让一个超荷变成标量}；由此得到的 $\mathbb{CP}^1$ 族共享同一配分函数，内禀参数是 $\Psi$；二维化后 $\mathbb{CP}^1_g$ 嵌入更大族，广义复几何把 A／B 两型统一起来\pC\anchor{src p.18–37, p.58–70}。
